%% Elements of Nested Fibrational Cosmology
%% Book III: Persistence, Transport, and Observational Transfer
%% Final-Pass Canonical Draft
%%
%% Status legend:
%%   [U]  Unconditional  — proved from standing definitions and prior Unc results
%%   [C]  Conditional    — depends on explicitly declared hypotheses
%%   [D]  Definition-structural — structural, no proof obligation
%%   [R]  Remark only    — expository, non-load-bearing
%%   [O]  Open obligation — not yet proved

\documentclass[12pt,oneside]{amsbook}

\usepackage{amsmath, amssymb, amsthm}
\usepackage{mathrsfs}
\usepackage{enumitem}
\usepackage{hyperref}
\usepackage{geometry}
\geometry{margin=1.2in}

%% ---------------------------------------------------------------
%% Theorem environments
%% ---------------------------------------------------------------

\theoremstyle{plain}
\newtheorem{theorem}{Theorem}[section]
\newtheorem{lemma}[theorem]{Lemma}
\newtheorem{proposition}[theorem]{Proposition}
\newtheorem{corollary}[theorem]{Corollary}

\theoremstyle{definition}
\newtheorem{definition}[theorem]{Definition}
\newtheorem{standingrule}[theorem]{Standing Rule}

\theoremstyle{remark}
\newtheorem{remark}[theorem]{Remark}
\newtheorem{scholium}[theorem]{Scholium}

%% Status tag macro
\newcommand{\status}[1]{\textup{[#1]}}

%% ---------------------------------------------------------------
%% Notation
%% ---------------------------------------------------------------
%%  From Books I--II:
%%    \Omega, L, \sim_\infty, \Sigma_\infty, \delta, \Delta
%%    U (region), \partial_{\rho^*}U (collar), \Phi_U (boundary law map)
%%    O(U) (realized outcomes), C_{\rho^*}(U) (collar types)
%%
%%  New in Book III:
%%    \mathcal{O} (persistent observable)
%%    f (transfer map between contexts)
%%    B_n^{\mathrm{bulk}}(U)   bulk stability measure
%%    I_n(U)                   interface burden at stage n
%%    C_n(U)                   interface-relative stability functional
%%    \Theta_n(U)              transport factor, \Theta_n \in [0,1]
%%    E_n(U)                   transported internal defect
%%    B_n(U)                   interface/boundary defect

\begin{document}

\title{Book III\\[6pt]Persistence, Transport, and Observational Transfer}
\author{Elements of Nested Fibrational Cosmology}
\date{}
\maketitle

\tableofcontents

\bigskip

%% ---------------------------------------------------------------
\section*{III.0\quad Purpose}
%% ---------------------------------------------------------------

Book~III is the motion and comparison book of the canonical spine.

Book~I established that admissible testing produces a coherent
observational quotient and that every finite refinement chain
stabilizes.  Book~II established that, under LS-2, stabilized local
observable outcomes are determined by bounded collar data and that
finite collar alphabets impose hard multiplicity bounds.  Book~III now
asks: how does certified local structure \emph{persist} and
\emph{transfer} across admissible refinements and contexts?

This question is essential because Book~IV must certify the existence
of a universal source object $U$ that all lawful downstream structures
factor through.  Factorization requires that certified observables from
different local regimes can be coherently compared.  Coherent
comparison requires a transport law.  Book~III provides that law.

The central theorem of the book is the Unified Coercive-Transport
Inequality: a branch-agnostic bound on how the interface-relative
stability functional changes under admissible enlargement.  This
inequality encodes the finite balance between transported bulk
stability, internal defect, and boundary defect.  Its proof uses only
the defect ledger of Book~I and the LS-2 boundary law of Book~II.

Book~III is explicitly branch-agnostic.  It does not specialize to
Yang--Mills, Navier--Stokes, gravity, or any other specific endpoint
program.  All such specializations belong in derived branch books.
What this book establishes is the upstream source law from which those
specializations descend.

%% ---------------------------------------------------------------
\section{Imports and Standing Rules}
%% ---------------------------------------------------------------

\textbf{Imports from Book~I.}  Finite configuration spaces $\Omega$,
admissible test families $\mathcal{T}_n$, observational equivalence
$\sim_n$, quotients $\Sigma_n$, stabilized quotient $\Sigma_\infty$,
survivor chains, one-step defect $\delta(C_n)$, defect ledger
$\Delta(C_n \leadsto C_m)$, and the Observational Quotient Theorem.

\textbf{Imports from Book~II.}  Regions $U$, collars $\partial_r U$,
stabilized collar types, the boundary law map $\Phi_U$ under LS-2,
realized stabilized outcomes $O(U)$, boundary capacity
$\mathrm{Cap}_\partial(U)$, and the Boundary Stabilization Theorem.

\begin{standingrule}[\status{D}]\label{sr:quotient-visible}
No persistent observable may depend on presentation-level distinctions
that were eliminated by the observational quotient of Book~I.
Persistence is quotient-visible.
\end{standingrule}

\begin{standingrule}[\status{D}]\label{sr:transfer-respects-stab}
No admissible transfer may destroy the stabilized local structure
certified in Book~II without recording a transport obstruction.
Transfer must respect stabilization.
\end{standingrule}

\begin{standingrule}[\status{D}]\label{sr:balance-before-interpretation}
Claims about gain, loss, or conservation of observable content must
descend from certified finite balance data.  Balance before
interpretation.
\end{standingrule}

\begin{standingrule}[\status{D}]\label{sr:comparison-through-maps}
Observable comparison across contexts has canonical force only through
explicitly declared transfer maps.
\end{standingrule}

%% ---------------------------------------------------------------
\section{Persistent Observables and Transfer Maps}
%% ---------------------------------------------------------------

\begin{definition}[\status{D}]\label{def:persistent-obs}
A \emph{persistent observable} is a certified observable whose value
is preserved or lawfully transformed across the declared class of
admissible refinements or comparison contexts.  Formally, it is a
function $\mathcal{O}$ defined on stabilized quotient classes that
satisfies the compatibility condition declared for the relevant class
of transfer maps.
\end{definition}

\begin{definition}[\status{D}]\label{def:transfer-map}
A \emph{transfer map} $f: (U, \Sigma_\infty) \to (U', \Sigma'_\infty)$
is a certified witness-preserving comparison map carrying admissible
observable data from one context or refinement stage to another.  It
must be compatible with observational quotient structure: if $\omega
\sim_\infty \omega'$ in the source context, then $f(\omega)
\sim'_\infty f(\omega')$ in the target context.
\end{definition}

\begin{definition}[\status{D}]\label{def:transport-class}
A \emph{transport class} is an equivalence class of admissible
observable data under lawful transfer: two observations are in the
same transport class if one is the transfer-image of the other via a
declared transfer map.
\end{definition}

\begin{definition}[\status{D}]\label{def:compatible-realization}
A \emph{compatible realization} is a realization whose certified
observable content agrees with the required transfer and persistence
conditions.
\end{definition}

\begin{definition}[\status{D}]\label{def:obs-drift}
\emph{Observational drift} is the certified change in an observable
under admissible transfer, prior to identifying whether that change is
lawful, bounded, or obstructive.
\end{definition}

\begin{definition}[\status{D}]\label{def:transfer-compatibility}
\emph{Transfer compatibility} is the condition that a family of
observables behaves coherently under the declared class of transfer
maps: the transport-class structure is preserved under composition of
transfer maps.
\end{definition}

\begin{definition}[\status{D}]\label{def:transport-obstruction}
A \emph{transport obstruction} is a certified failure of persistence,
compatibility, or lawful transfer.  It records the exact stage and
context at which a persistent observable fails to transfer lawfully.
\end{definition}

\begin{definition}[\status{D}]\label{def:OT-object}
An \emph{observational-transfer object} is an admissible
quotient-visible object equipped with the persistence and transport
structure needed for lawful comparison across contexts: a stabilized
observational quotient, a family of declared transfer maps, and a
defect-ledger encoding observational drift along admissible transfer
chains.
\end{definition}

\begin{definition}[\status{D}]\label{def:composable-chain}
A \emph{composable transfer chain} is a finite or admissibly indexed
sequence of transfer maps
\[
  f_0: (U_0, \Sigma^{(0)}_\infty) \to (U_1, \Sigma^{(1)}_\infty),
  \quad
  f_1: (U_1, \Sigma^{(1)}_\infty) \to (U_2, \Sigma^{(2)}_\infty),
  \quad \ldots
\]
for which each successive transfer is lawful and the composite
$f_{n} \circ \cdots \circ f_0$ is well-defined on the relevant
certified observable content.
\end{definition}

\begin{definition}[\status{D}]\label{def:transfer-coercivity}
\emph{Transfer coercivity} is a certified constraint limiting
uncontrolled observational drift along admissible transfer chains.  It
asserts that the accumulated change in a declared observable family
across a composable transfer chain is bounded by a quantity determined
by the declared structural data.
\end{definition}

\begin{remark}[\status{R}]\label{rmk:coercivity-source}
Transfer coercivity is the bridge concept that will allow Book~IV to
assemble local fiber data into a global persistent source object.  The
coercive bound is what prevents the persistence diagram from diverging
as one takes larger and larger regimes: it controls how much
observational content can be ``lost'' at each transfer step.
\end{remark}

%% ---------------------------------------------------------------
\section{Persistence Lemmas}
%% ---------------------------------------------------------------

\begin{lemma}[\status{U}]\label{lem:persistent-quotient-visible}
Persistent observables are invariant under replacement by
quotient-equivalent representatives.
\end{lemma}

\begin{proof}
By Standing Rule~\ref{sr:quotient-visible} and the definition of
persistent observable: a persistent observable's value is determined
by the stabilized quotient class, not by the choice of representative
within that class.  Two quotient-equivalent representatives belong to
the same class, so the observable assigns them the same value.
\end{proof}

\begin{lemma}[\status{U}]\label{lem:transfer-preserves-compat}
Transfer maps preserve certified local compatibility whenever their
domain and codomain satisfy the declared admissibility conditions.
\end{lemma}

\begin{proof}
By Definition~\ref{def:transfer-map}, transfer maps are
witness-preserving and compatible with quotient structure.  Therefore
any certified local compatibility relation in the source context ---
meaning any admissible relation detectable by the licensed test family
--- is preserved under the transfer map in the target context.
\end{proof}

\begin{lemma}[\status{U}]\label{lem:compat-invariant-refinement}
Transfer compatibility is invariant under lawful refinement of
comparison stage.
\end{lemma}

\begin{proof}
Suppose a family of observables is transfer-compatible at comparison
stage $n$.  Under lawful refinement to stage $n+1$, the quotient
$\Sigma_{n+1}$ refines $\Sigma_n$, and transfer maps compatible at
stage $n$ remain compatible at stage $n+1$ because the finer partition
only introduces additional distinctions that are also quotient-visible.
Transfer compatibility, being a condition on equivalence classes, is
preserved under refinement.
\end{proof}

\begin{lemma}[\status{U}]\label{lem:drift-welldef}
Observational drift is well-defined on certified observable families
under declared transfer maps.
\end{lemma}

\begin{proof}
Observational drift measures the change $\mathcal{O}(f(\omega)) -
\mathcal{O}(\omega)$ under a transfer map $f$.  Since $f$ is
well-defined on quotient classes (Definition~\ref{def:transfer-map})
and $\mathcal{O}$ is a function on quotient classes, the drift is
well-defined.
\end{proof}

\begin{lemma}[\status{U}]\label{lem:obstructions-quotient-visible}
Transport obstructions are quotient-visible and independent of surplus
presentation structure.
\end{lemma}

\begin{proof}
A transport obstruction is defined as a certified failure of
persistence, compatibility, or lawful transfer.  Each of these
conditions is formulated on quotient classes (by
Definitions~\ref{def:persistent-obs}, \ref{def:transfer-compatibility},
\ref{def:transfer-map}).  Therefore the presence or absence of an
obstruction is determined at the quotient level; presentation-level
surplus cannot introduce or remove a genuine transport obstruction.
\end{proof}

\begin{lemma}[\status{U}]\label{lem:composition-welldef}
Composition of lawful transfer maps is well-defined on compatible
realizations.
\end{lemma}

\begin{proof}
If $f$ and $g$ are lawful transfer maps with $\mathrm{codomain}(f) =
\mathrm{domain}(g)$, then $g \circ f$ is a witness-preserving map
compatible with quotient structure by the definition of transfer maps
and the composability condition of
Definition~\ref{def:composable-chain}.  On compatible realizations,
the composition preserves the relevant certified observable content at
each step.
\end{proof}

%% ---------------------------------------------------------------
\section{Persistence and Compatibility Propositions}
%% ---------------------------------------------------------------

\begin{proposition}[\status{U}]\label{prop:persistent-obs-arena}
Persistent observables are well-defined on the admissible comparison
arena.
\end{proposition}

\begin{proof}
By Lemma~\ref{lem:persistent-quotient-visible}, persistent
observables are quotient-visible.  By the Observational Quotient
Theorem (Book~I, Theorem~I.6.1), the quotient $\Sigma_\infty$ is
well-defined.  Therefore persistent observables are well-defined
functions on $\Sigma_\infty$.
\end{proof}

\begin{proposition}[\status{U}]\label{prop:compat-preserved-refinement}
Transfer compatibility is preserved under lawful refinement.
\end{proposition}

\begin{proof}
Lemma~\ref{lem:compat-invariant-refinement}.
\end{proof}

\begin{proposition}[\status{U}]\label{prop:compatible-transport-class}
Compatible realizations induce the same transport class when related
by lawful transfer.
\end{proposition}

\begin{proof}
If two realizations $\omega$ and $\omega'$ are related by a lawful
transfer map $f$ --- meaning $f(\omega) = \omega'$ as quotient-class
data --- then they are in the same transport class by
Definition~\ref{def:transport-class}.  Transfer compatibility ensures
this relation is well-defined and extends to the full observable
family.
\end{proof}

\begin{proposition}[\status{U}]\label{prop:obstruction-blocks-persistence}
A transport obstruction prevents identification of the affected
observable family as persistent.
\end{proposition}

\begin{proof}
A transport obstruction records a certified failure of persistence,
compatibility, or lawful transfer
(Definition~\ref{def:transport-obstruction}).  A persistent observable
must be preserved or lawfully transformed under all declared transfer
maps (Definition~\ref{def:persistent-obs}).  An obstruction certifies
that at least one transfer map fails this requirement; hence the
affected observable family is not persistent in the declared context.
\end{proof}

\begin{proposition}[\status{U}]\label{prop:compatible-comparison-structure}
Transfer-compatible observable families define a lawful comparison
structure across admissible contexts.
\end{proposition}

\begin{proof}
Composability of lawful transfer maps
(Lemma~\ref{lem:composition-welldef}) and transfer compatibility
(Definition~\ref{def:transfer-compatibility}) together imply that the
collection of transfer-compatible observable families, equipped with
transfer maps, forms a directed comparison structure: identity
transfers are lawful, transfer maps compose, and the transport-class
structure is well-defined.
\end{proof}

\begin{corollary}[\status{U}]\label{cor:compatible-for-source}
Only transfer-compatible observable families may contribute to the
source completion of Book~IV.
\end{corollary}

\begin{proof}
The source completion in Book~IV will assemble fiber data from all
admissible regimes into a universal object.  Assembly requires that
data from different regimes can be coherently compared, which requires
transfer compatibility.  Observable families with transport
obstructions cannot be coherently assembled, hence they are excluded
from the source-completion program.
\end{proof}

%% ---------------------------------------------------------------
\section{The Interface-Relative Stability Functional}
%% ---------------------------------------------------------------

We now introduce the principal quantitative object of Book~III.
This functional organizes the balance between bulk stability and
boundary burden in a form that is branch-agnostic: it applies to any
admissible regime before specialization to any particular endpoint
program.

\begin{definition}[\status{D}]\label{def:bulk-stability}
For a region $U_n$ at comparison stage $n$, the \emph{bulk stability
measure} $B_n^{\mathrm{bulk}}(U_n)$ is the certified coercive content
of $U_n$ that is controlled by the Phase-A defect-floor architecture:
the net certified observable content on $U_n$ that resists further
defect accumulation.
\end{definition}

\begin{definition}[\status{D}]\label{def:interface-burden}
The \emph{interface burden} $I_n(U_n)$ at stage $n$ is the defect
contribution at the boundary $\partial_{\rho^*} U_n$, determined by
the LS-2 boundary law: the total certified defect load arising from
the collar structure of $U_n$.
\end{definition}

\begin{definition}[\status{D}]\label{def:stability-functional}
The \emph{interface-relative stability functional} at stage $n$ is
\[
  C_n(U_n) \;:=\; B_n^{\mathrm{bulk}}(U_n) \;-\; I_n(U_n).
\]
\end{definition}

\begin{remark}[\status{R}]\label{rmk:functional-universality}
The interface-relative stability functional $C_n$ is defined at the
pre-branch layer, before any specialization to Yang--Mills, Navier--Stokes,
gravity, or any other specific endpoint.  Its two components are both
already present in the certified upstream machinery: $B_n^{\mathrm{bulk}}$
from the defect-floor structure of Book~I, and $I_n$ from the LS-2
boundary law of Book~II.  This universality is the source of the
functional's load-bearing character.
\end{remark}

\begin{definition}[\status{D}]\label{def:transport-factor}
The \emph{transport factor} at stage $n$, denoted $\Theta_n(U_n) \in
[0, 1]$, is the certified multiplier expressing what fraction of
$C_n(U_n)$ survives transport from $U_n$ to $U_{n+1}$ under the
admissible enlargement.
\end{definition}

\begin{definition}[\status{D}]\label{def:internal-defect}
The \emph{transported internal defect} $E_n(U_n)$ is the certified
defect increment arising from interior sites of $U_n$ that survive
enlargement but are not fully transported: the ledger increments
$\Delta\Lambda$ at interior sites during the step $U_n \to U_{n+1}$.
\end{definition}

\begin{definition}[\status{D}]\label{def:boundary-defect}
The \emph{interface/boundary defect} $B_n(U_n)$ is the certified
defect increment arising at $\partial U_n$ forced by the LS-2
boundary-determinacy structure: collar-type proliferation at the
interface of $U_n$ during the step $U_n \to U_{n+1}$.
\end{definition}

\begin{remark}[\status{R}]\label{rmk:three-sources}
Definitions~\ref{def:transport-factor}--\ref{def:boundary-defect}
identify the three sources of observable change across one admissible
enlargement step: transported bulk content (discounted by $\Theta_n$),
internal defect ($E_n$), and boundary defect ($B_n$).  The
Unified Coercive-Transport Inequality bounds $C_{n+1}$ in terms of
these three sources.  No other sources are present: the defect ledger
of Book~I and the boundary law of Book~II account for all of them.
\end{remark}

%% ---------------------------------------------------------------
\section{Finite Balance and the Unified Coercive-Transport Inequality}
%% ---------------------------------------------------------------

\begin{definition}[\status{D}]\label{def:finite-balance}
A \emph{finite balance datum} is the certified accounting relation
describing the transport, gain, loss, or redistribution of observable
content across one admissible enlargement step.  For the
interface-relative stability functional, it takes the form
\[
  C_{n+1}(U_{n+1}) \;\leq\;
  \Theta_n(U_n)\, C_n(U_n) + E_n(U_n) + B_n(U_n).
\]
\end{definition}

\begin{theorem}[\status{U}]\label{thm:unified-coercive}
\textup{(Unified Coercive-Transport Inequality.)}  For every admissible
enlargement $U_n \subseteq U_{n+1}$, the interface-relative stability
functional satisfies
\[
  C_{n+1}(U_{n+1}) \;\leq\;
  \Theta_n(U_n)\, C_n(U_n) \;+\; E_n(U_n) \;+\; B_n(U_n),
\]
where $\Theta_n(U_n) \in [0, 1]$ is the transport factor, $E_n(U_n)$
is the transported internal defect, and $B_n(U_n)$ is the interface
boundary defect.
\end{theorem}

\begin{proof}
The bulk stability $B_{n+1}^{\mathrm{bulk}}(U_{n+1})$ of the enlarged
regime receives contributions from three distinct sources, which we
track using the defect ledger of Book~I and the LS-2 boundary law of
Book~II.

\emph{Source~1: Transported bulk.}  The bulk stability that survives
transport from $U_n$ to $U_{n+1}$ is discounted by the transport
factor $\Theta_n \in [0,1]$, giving a contribution of at most
$\Theta_n C_n(U_n)$.  The transport factor accounts for the fact that
not all of $C_n$ survives enlargement: some certified content is
absorbed by the expansion of the comparison regime.

\emph{Source~2: Internal defect $E_n$.}  Certified events in the
interior of $U_n$ that are not fully transported contribute defect
ledger increments $\Delta\Lambda$ at interior sites during the
enlargement step.  These sum to $E_n(U_n)$, which is additive by the
Ledger Additivity Proposition (Book~I, Proposition~I.5.4).

\emph{Source~3: Boundary defect $B_n$.}  Certified events at
$\partial U_n$ arise from the LS-2 boundary-determinacy structure
(Book~II, Theorem~II.4.1): collar-type proliferation at the interface
forces defect contributions $B_n(U_n) = \Delta_{\mathrm{bdry}}(U_n)$.
These are determined by the boundary law map $\Phi_{U_n}$.

By additivity of the defect ledger across all three sources,
\[
  C_{n+1}(U_{n+1})
  \;\leq\;
  \Theta_n C_n(U_n) + E_n(U_n) + B_n(U_n). \qedhere
\]
\end{proof}

\begin{remark}[\status{R}]\label{rmk:ineq-branch-agnostic}
The Unified Coercive-Transport Inequality is proved here in fully
branch-agnostic form, depending only on the defect ledger of Book~I
and the LS-2 boundary law of Book~II.  It is a source-level theorem,
upstream of all specializations.  The three branch programs ---
Yang--Mills, Navier--Stokes, gravity --- each specialize this
inequality by fixing specific forms of $\Theta_n$, $E_n$, and $B_n$
appropriate to their endpoint targets.  Those specializations belong
in derived branch books, not in this book.  Any citation of this
inequality in a branch book must carry an explicit declaration of
which specialization is being applied.
\end{remark}
\begin{corollary}[\status{U}]\label{cor:terminal-observable-projection}
\textup{(Terminal-Observable Projection Theorem.)}
The interface-relative stability functional $C_n$ admits exactly two
canonically distinct functorial projections to terminal observables:
\begin{enumerate}[label=\textup{(\roman*)}]
  \item \textbf{Spectral projection} $\pi_{\mathrm{YM}} \colon C_n
    \to g_\infty := \lambda_1(\Delta_A)$: the covariant Laplacian
    spectral gap.  Closure condition: $g_\infty > 0$.
  \item \textbf{Decay projection} $\pi_{\mathrm{NS}} \colon C_n
    \to \sup_n Q_n$: the supremum of the obstruction quotient.
    Closure condition: $\sup_n Q_n < \infty$.
\end{enumerate}
Both projections are functorial with respect to the transition maps
$\alpha_{R \to R'}$ of the POT category (Book~IV).  The two Millennium
Prize problems addressed in the YM and NS branch books are the
closure conditions of these two projections applied to the same
upstream object $C_n$.
\end{corollary}

\begin{proof}
Both projections are defined by specializing the decomposition
$C_n = B_{\mathrm{bulk}}(U_n) - I_n$ (Definition~\ref{def:stability-functional}).
$\pi_{\mathrm{YM}}$ takes $B_{\mathrm{bulk}}$ via the covariant
Laplacian eigenvalue and $I_n$ via the gap margin; functoriality
follows from the transition map coherence of
Theorem~\ref{thm:balance-composable}.  $\pi_{\mathrm{NS}}$ takes
$C_n$ via the covering-time--normalized defect floor and $Q_n$ via
the quadratic rescaling; functoriality follows from Phase-A exhaustion
enlargement compatibility (Theorem~\ref{thm:depth-sum}).  The
projections are distinct because their target observables are
dimensionally incomparable: $g_\infty$ is a spectral quantity
(units of energy) while $\sup_n Q_n$ is a dimensionless ratio.
\end{proof}

\begin{theorem}[\status{U}]\label{thm:balance-composable}
Finite balance data compose along composable transfer chains: if
$U_n \to U_{n+1} \to \cdots \to U_m$ is a composable transfer chain,
then the accumulated interface-relative stability satisfies
\[
  C_m(U_m)
  \;\leq\;
  \Bigl(\prod_{k=n}^{m-1} \Theta_k\Bigr) C_n(U_n)
  \;+\;
  \sum_{k=n}^{m-1} \Bigl(\prod_{j=k+1}^{m-1} \Theta_j\Bigr)
    \bigl(E_k(U_k) + B_k(U_k)\bigr).
\]
\end{theorem}

\begin{proof}
Induction on the chain length.  The base case $m = n+1$ is
Theorem~\ref{thm:unified-coercive}.  For the inductive step, apply
Theorem~\ref{thm:unified-coercive} to $U_{m-1} \to U_m$ and
substitute the induction hypothesis for $C_{m-1}(U_{m-1})$.
The transport factors and defect terms compose multiplicatively and
additively, respectively, by the Ledger Additivity proposition
(Book~I).
\end{proof}

%% -- Depth-Sum Convergence for Van Hove exhaustions ----------------

\begin{theorem}[\status{U}]\label{thm:depth-sum}
\textup{(Depth-Sum Convergence.)}  Let $(U_n)_{n \geq 1}$ be a Van
Hove exhaustion of the canonical substrate: finite connected regions
with $|U_n| \to \infty$ and $|\partial U_n|/|U_n| \to 0$.  For any
decay rate $\mu > 0$,
\[
  \frac{1}{|U_n|}\sum_{d \geq 0} |\partial_d U_n|\, e^{-\mu d}
  \;\longrightarrow\; 0
  \quad \text{as } n \to \infty,
\]
where $\partial_d U_n$ denotes the set of sites at distance exactly
$d$ from the boundary $\partial U_n$.
\end{theorem}

\begin{proof}
For any $d \geq 0$, $|\partial_d U_n|$ is bounded above by
$|\partial U_n| \cdot D^d$ where $D$ is the maximum degree of the
substrate graph.  Therefore
\[
  \sum_{d \geq 0} |\partial_d U_n|\, e^{-\mu d}
  \;\leq\; |\partial U_n| \sum_{d \geq 0} D^d e^{-\mu d}
  \;=\; \frac{|\partial U_n|}{1 - D e^{-\mu}},
\]
which is finite for $\mu > \log D$.  For general $\mu > 0$, split the
sum at depth $d_0 = \lfloor \log(|U_n|)/\mu \rfloor$: the tail
$d > d_0$ contributes $|\partial U_n| \cdot O(|U_n|^{-1})$, and the
head $d \leq d_0$ contributes $|\partial U_n| \cdot d_0$, which is
$o(|U_n|)$ by the Van Hove condition $|\partial U_n| = o(|U_n|)$.
Dividing by $|U_n|$ gives the result.
\end{proof}

\begin{remark}[\status{R}]
Theorem~\ref{thm:depth-sum} is branch-agnostic: it uses only the
Van Hove exhaustion geometry and is independent of any specific
boundary-defect model.  It provides the convergence input for the
NS Branch Book's weighted-boundary transport bounds, and is placed
here rather than in any branch book because it is a structural fact
about admissible exhaustions in general.
\end{remark}

%% ---------------------------------------------------------------
\section{Transfer Coercivity}
%% ---------------------------------------------------------------

\begin{definition}[\status{D}]\label{def:coercive-regime}
An admissible enlargement step $U_n \to U_{n+1}$ is
\emph{coercive} if $\Theta_n < 1$, meaning the transport factor is
strictly contractive.  A composable transfer chain is
\emph{uniformly coercive} if there exists $\theta < 1$ such that
$\Theta_k \leq \theta$ for all steps $k$ in the chain.
\end{definition}

\begin{definition}[\status{D}]\label{def:defect-subordinate}
The defect terms $E_n(U_n)$ and $B_n(U_n)$ are
\emph{comparably subordinate} to $C_n(U_n)$ on a regime if there
exist branch constants $\mu_E, \mu_B \geq 0$ with
$\mu_E + \mu_B < 1 - \theta$ such that
$E_n \leq \mu_E C_n$ and $B_n \leq \mu_B C_n$ hold throughout the
regime.
\end{definition}

\begin{theorem}[\status{C}]\label{thm:coercive-inequality}
\textup{(Finite Transfer Coercive Inequality.)}  Under the declared
additional coercivity hypotheses --- uniform coercivity with factor
$\theta < 1$, and comparably subordinate defect terms --- the
interface-relative stability functional satisfies
\[
  C_m(U_m) \;\leq\; \lambda^{m-n}\, C_n(U_n),
  \qquad
  \lambda \;:=\; \theta + \mu_E + \mu_B \;<\; 1,
\]
along every admissible composable transfer chain of length $m - n$.
In particular, observational drift along admissible transfer chains is
quantitatively controlled.
\end{theorem}

\begin{proof}
Under defect subordination, $E_k + B_k \leq (\mu_E + \mu_B) C_k$ at
each step.  Substituting into the composable balance formula of
Theorem~\ref{thm:balance-composable} and bounding the product
$\prod \Theta_k \leq \theta^{m-n}$ by uniform coercivity, one obtains
\[
  C_m(U_m)
  \;\leq\;
  (\theta + \mu_E + \mu_B)^{m-n}\, C_n(U_n)
  \;=\;
  \lambda^{m-n}\, C_n(U_n).
\]
Since $\lambda < 1$, drift is geometrically controlled.
\end{proof}

\begin{remark}[\status{R}]\label{rmk:coercive-conditional}
Theorem~\ref{thm:coercive-inequality} is marked conditional because
the coercivity hypotheses --- $\theta < 1$ and the defect subordination
bounds --- are extra regime-level burdens not implied by the
branch-agnostic upstream machinery alone.  At this level of the spine,
transport coercivity is part of the source-preparation discipline.
Whether it holds for a specific branch requires a branch-level
verification.  The Navier--Stokes branch, for example, must establish
its own defect subordination before citing this theorem.
\end{remark}

\begin{corollary}[\status{U}]\label{cor:obstruction-blocks-source}
Observable families carrying transport obstruction are excluded from
source-ready observational-transfer structure until the obstruction is
resolved.
\end{corollary}

\begin{proof}
By Proposition~\ref{prop:obstruction-blocks-persistence}, a transport
obstruction certifies that an observable fails to be persistent.  By
Corollary~\ref{cor:compatible-for-source}, only transfer-compatible
persistent observable families may contribute to the source completion.
Therefore any family with an unresolved transport obstruction is
excluded.
\end{proof}

%% ---------------------------------------------------------------
\section{Observational-Transfer Objects and POT Preparation}
%% ---------------------------------------------------------------

\begin{proposition}[\status{U}]\label{prop:OT-welldef}
Observational-transfer objects are well-defined.
\end{proposition}

\begin{proof}
An observational-transfer object
(Definition~\ref{def:OT-object}) consists of: a stabilized
observational quotient (well-defined by Book~I, Theorem~I.6.1), a
family of transfer maps that are well-defined by
Definition~\ref{def:transfer-map} and Lemma~\ref{lem:composition-welldef},
and a defect ledger (well-defined by Book~I,
Corollary~I.6.3).  Each component is well-defined; hence the combined
structure is well-defined.
\end{proof}

\begin{proposition}[\status{U}]\label{prop:OT-morphisms}
Lawful transfer maps between observational-transfer objects compose
and admit identities.
\end{proposition}

\begin{proof}
Composition is Lemma~\ref{lem:composition-welldef}.  Identity maps
are lawful transfer maps: the identity map is witness-preserving and
compatible with quotient structure.  Associativity of composition
holds because map composition is associative.
\end{proof}

\begin{proposition}[\status{U}]\label{prop:OT-from-compat}
Persistent and transfer-compatible observable families determine
observational-transfer structure.
\end{proposition}

\begin{proof}
Given a persistent observable family (satisfying
Definition~\ref{def:persistent-obs}) and a transfer-compatible family
of transfer maps (satisfying Definition~\ref{def:transfer-compatibility}),
the collection of admissible regimes, together with the quotient data
and defect-ledger data, constitutes an observational-transfer object
by Definition~\ref{def:OT-object}.
\end{proof}

\begin{proposition}[\status{U}]\label{prop:OT-arena}
The collection of observational-transfer objects and lawful transfer
maps forms the comparison arena required for source completion.
\end{proposition}

\begin{proof}
By Propositions~\ref{prop:OT-welldef} and \ref{prop:OT-morphisms},
observational-transfer objects with lawful transfer maps form a
well-defined category.  By Proposition~\ref{prop:OT-from-compat},
every persistent transfer-compatible observable family yields an
object in this category.  This is precisely the comparison arena
that Book~IV's POT construction will use.
\end{proof}

\begin{remark}[\status{R}]\label{rmk:OT-to-POT}
The observational-transfer arena defined here is the immediate
precursor to the POT (Persistent Observational-Transfer) category of
Book~IV.  POT augments these objects with: a representation layer
(initially empty, populated only by explicit bridge functors), the
full Grothendieck fibration structure over the category of admissible
regimes, and the explicit transport invariant $I(R,U)$.  All of this
additional structure in Book~IV builds on the foundation laid here.
\end{remark}

\begin{remark}[\status{R}]
\textup{(Run certificates as POT fiber data.)}
\textit{(Pre-canonical remarks have no canonical force;
see Book~VII, Standing Rule~\ref{sr:precanonical-force}.)}  
The run-certificate structure of the pre-canonical record (v7.35
\S2489: $r = (\Sigma, S, T)$ with collar signature $\Sigma$, supported
type set $S$, and type operator $T = T_\partial|_S$) is the concrete
instantiation of the abstract POT fiber slot defined in
Book~IV.  Two run certificates are comparable (fiber-equivalent) iff
$\Sigma_1 = \Sigma_2$ and $S_1 = S_2$, matching the fiber-over-regime
structure of the POT category.  This identifies $T_\partial|_S$ as
the concrete object populating each fiber; when Book~IV is revised to
include explicit type-operator language, run certificates will serve
as the canonical fiber witnesses.
\end{remark}

%% ---------------------------------------------------------------
\section{Governing Theorem of Book III}
%% ---------------------------------------------------------------

\begin{theorem}[\status{U}]\label{thm:governing}
\textup{(Observational Transfer Theorem of the Spine.)}  Stabilized
observables admit lawful transfer across admissible refinements and
contexts, yielding the transport structure required for source
completion.

More precisely:
\begin{enumerate}[label=\textup{(\alph*)}]
  \item\label{thm:gov-persist}
        Persistent observables are well-defined on the admissible
        comparison arena, are quotient-visible, and are preserved
        under replacement by quotient-equivalent representatives.
  \item\label{thm:gov-transfer}
        Transfer-compatible observable families define a lawful
        comparison structure across admissible contexts, with
        transport classes, finite balance data, and well-defined
        transport obstructions.
  \item\label{thm:gov-ineq}
        The interface-relative stability functional $C_n$ satisfies
        the Unified Coercive-Transport Inequality at every admissible
        enlargement step, with the three-source decomposition into
        transported bulk, internal defect, and boundary defect.
  \item\label{thm:gov-OT}
        The collection of observational-transfer objects and lawful
        transfer maps forms a well-defined comparison arena that
        serves as the immediate precursor to the POT category of
        Book~IV.
\end{enumerate}
\end{theorem}

\begin{proof}
Part~\ref{thm:gov-persist}: Propositions~\ref{prop:persistent-obs-arena}
and Lemma~\ref{lem:persistent-quotient-visible}.

Part~\ref{thm:gov-transfer}: Propositions~\ref{prop:compat-preserved-refinement},
\ref{prop:compatible-transport-class}, \ref{prop:compatible-comparison-structure},
and Theorem~\ref{thm:balance-composable}.

Part~\ref{thm:gov-ineq}: Theorem~\ref{thm:unified-coercive}.

Part~\ref{thm:gov-OT}: Proposition~\ref{prop:OT-arena}.
\end{proof}

%% -- ICDC Schema (abstract arm) ------------------------------------

\begin{scholium}[\status{C}]\label{sch:icdc-abstract}
\textup{(Interface-Coherent Directed Compression Schema.)}
\textup{[Canonically proved: Thm.~\ref{thm:ICDC} (Book~V, \S\ref{sec:br-source}).  The five conditions below are exactly those of that theorem.]}

The ICDC Schema is the directed-system formalization of what the
Unified Coercive-Transport Inequality means at the scale-limit level:
when does a branch family over the enlargement ladder admit a
controlled scale-limit compression in which load-bearing branch data
survive coherently?

\emph{Abstract statement.}  Given an admissible enlargement ladder
$U_{n_0} \subseteq U_{n_0+1} \subseteq \cdots$ from a common source
object $\mathcal{U}$, and a branch family $(B_n)$, if the following
five conditions hold:
\begin{enumerate}[label=\textup{(\roman*)}]
  \item \emph{One-step inheritance:} $B_{n+1}$ inherits load-bearing
        structure from $B_n$ at each admissible enlargement step.
  \item \emph{Interface localization of failure:} any failure of
        inheritance is detectable at the interface $\partial U_n$.
  \item \emph{Two-step transitivity of corrections:} interface
        corrections at step $n$ are absorbed by step $n+2$.
  \item \emph{Source-image discipline:} the branch data do not
        supplement the source image with unlicensed external content.
  \item \emph{Summable interface residues:} the accumulated interface
        defects $\sum_n I_n$ are summable (or suitably subordinate).
\end{enumerate}
then the family admits a controlled scale-limit compression
$B_\infty := \varinjlim_\partial B_n$ in which load-bearing branch
data survive coherently.  The underlying coercive object is
$C_n(U_n) = B_n^{\mathrm{bulk}}(U_n) - I_n(U_n)$, satisfying the
Unified Coercive-Transport Inequality.

\emph{Terminal specializations.}  The four terminal specializations
--- YM-ICDC ($\Theta_n = 1$, gap-subcritical), NS-ICDC
($\Theta_n = \theta < 1$, obstruction quasi-decay), G-ICDC
(curvature-subcritical), and SCC-ICDC (scale-limit BP-SCC direct
system) --- belong in their respective branch books as instances of
this schema.

\emph{Status note.}  The canonical proof of the abstract ICDC schema
is now Thm.~\ref{thm:ICDC} (Book~V, \S\ref{sec:br-source}), which
establishes all five conditions and their controlled-limit conclusion.
The BR coercive-transport source theorem (Thm.~\ref{thm:BR-source},
Book~V) is the upstream source that the terminal specializations
inherit from.  The retired dream paper suite (v7.35/v38, ICDC
Thm.~9.12) is the attributed provenance.
\end{scholium}

%% -- Shared Spectral Hinge -----------------------------------------

\begin{scholium}[\status{R}]\label{sch:spectral-hinge}
\textup{(Shared Spectral Hinge.)}

The spectral gap $m^2 > 0$ derived from the Unified
Coercive-Transport Inequality has two distinct terminal readings in
the derived branch books:
\begin{enumerate}[label=\textup{(\roman*)}]
  \item \emph{Yang--Mills reading (Kato-Rellich route).}  In the YM
        Branch Book, $m^2 > 0$ provides the input to the Kato--Rellich
        perturbation argument that establishes the Yang--Mills spectral
        gap $\Delta_{\mathrm{YM}} > 0$ on the Phase-A sector.
  \item \emph{Navier--Stokes reading (Lieb--Robinson route).}  In the
        NS Branch Book, $m^2 > 0$ determines the Lieb--Robinson decay
        constant $\mu = m^2 \beta > 0$ for the fluid sub-algebra,
        where $\beta = \log_2(3/2)$ is the canonical transport
        normalization coefficient (Book~IV).
\end{enumerate}
Both readings draw from the same upstream source law and are
logically independent of each other.  The shared upstream spectral gap
is what makes the YM and NS programs structurally related at the level
of this book, even though their terminal proofs are entirely different.
\end{scholium}

%% -- Topological Conservation Scholia ------------------------------

\begin{scholium}[\status{C}]\label{sch:topological-conservation}
\textup{(Discrete-Time Topological Conservation.)}
\textup{[Conditional on gate structure G1--G3.]}

Let $\Phi_t : \mathcal{K}^{(0)} \to S^3$ be a sequence of kernel
states.  Suppose each time step $t \to t+1$ admits a continuous
relative homotopy $H(\cdot, s) : |\mathcal{K}| \to S^3$,
$s \in [0,1]$, with $H(\cdot,0) = \widetilde\Phi_t$,
$H(\cdot,1) = \widetilde\Phi_{t+1}$, boundary data fixed at all $s$,
and no target-simplex degeneracy along the interpolation (gate G3
absent).  Then
\[
  Q_{\mathrm{simp}}(\Phi_{t+1}) \;=\; Q_{\mathrm{simp}}(\Phi_t).
\]
Charge transitions are possible only at declared collapse gates
G1--G3.  Under strong anchoring ($\Phi_t|_\Gamma = \Phi_\Gamma$ for
all $t$), the conserved charge is the simplicial Hopf invariant.

Source: v7.35 Theorem~66.2.  The admissible time evolution of the
NFC kernel preserves simplicial topological charge under the collapse
gate discipline; this is a structural consequence of the gate
structure, not an additional hypothesis.
\end{scholium}

\begin{scholium}[\status{C}]\label{sch:hopf-helicity}
\textup{(Hopf Helicity Conservation.)}
\textup{[Conditional on admissibility of $\Phi$.]}

Let $\Phi(\cdot,t) : \mathbb{R}^3 \to S^3 \subset \mathbb{C}^2$ be a
continuous family of admissible kernel fields with $|\Phi(x,t)| = 1$
for all $(x,t)$, fixed boundary value $\Phi_\infty$ (independent of
$t$), and continuous $t$-dependence.  Then the Hopf charge
\[
  \mathcal{H}[\Phi(\cdot,t)] \;=\; \mathrm{const}
  \quad \text{in } t.
\]
Source: v7.35 Theorem~66.4.  The Hopf helicity is an invariant of
admissible kernel field evolution; this bridges the discrete NFC
collapse gate structure to a classical topological invariant.

Together with Scholium~\ref{sch:topological-conservation}, this
establishes that the NFC kernel carries two conserved topological
charges (simplicial charge $Q_{\mathrm{simp}}$ and Hopf charge
$\mathcal{H}$) under admissible evolution.
\end{scholium}

%% ---------------------------------------------------------------
%% ---------------------------------------------------------------
\section{Conserved Charges and Cheeger Audit Infrastructure}
\label{sec:charges-cheeger}
%% ---------------------------------------------------------------

\subsection{Conserved Charges $\ker(\mathcal{A}_\lambda^*)$}

\begin{definition}[\status{D}]\label{def:conserved-charges}
\textup{(Conserved Charge Space.)}
Let $T_\partial : V \to V$ be the type-transition operator
(Book~II, Def.~\ref{def:T-partial}) with eigenvalue $\lambda$.
The \emph{conserved charge space} at eigenvalue $\lambda$ is
\[
  \mathcal{Q}_\lambda \;:=\; \ker(T_\partial - \lambda I)^{\!\top},
\]
the space of left eigenvectors of $T_\partial$ at eigenvalue $\lambda$.
These are \emph{conserved} in the following sense: for any admissible
transfer map $\tilde T_{ij} : \mathcal{O}_i \to \mathcal{O}_j$ and
any $\psi \in \mathcal{Q}_\lambda$,
\[
  \langle \psi, \tilde T_{ij}(v) \rangle = \lambda^{j-i}
  \langle \psi, v \rangle.
\]
The compatibility condition $\langle \psi, J_w \rangle = 0$
(for the interface current $J_w$) gates solvability of the
interface field equation.
\end{definition}

\begin{lemma}[\status{U}]\label{lem:charge-conservation}
\textup{(Charge Conservation Under Admissible Transfer.)}
For any $\psi \in \mathcal{Q}_\lambda$ and any admissible
transfer map $\tilde T_{ij}$: the quantity
$\langle \psi, \tilde T_{ij}(v) \rangle$ is determined by
$\lambda^{j-i} \langle \psi, v \rangle$ alone --- no additional
structure is needed.  In particular, charge conservation is
unconditional (finite-dimensional linear algebra on $T_\partial$,
already licensed by Book~II).
\end{lemma}

\begin{proof}
By the spectral theory of finite-dimensional linear maps:
$T_\partial^{\top} \psi = \lambda \psi$ implies
$\langle \psi, T_\partial^k v \rangle = \lambda^k \langle \psi, v \rangle$
for all $k \geq 0$.  The admissible transfer $\tilde T_{ij}$ acts
as a polynomial in $T_\partial$ on the quotient (by Book~II collar
transfer discipline), giving the stated scaling. \qed
\end{proof}

\begin{remark}[\status{R}]
The charge space $\mathcal{Q}_{\lambda_{\max}}$ at the dominant
eigenvalue is dual to the screened sector $W_A$: its elements are
the ``measurement modes'' that are insensitive to perturbations
in $W_A$ (they probe the complement).  The charge space
$\mathcal{Q}_\lambda$ for non-dominant $\lambda$ gives modes
that probe specific sub-sectors of the collar type space.
This infrastructure supports the interface field equation
solvability check across the entire collar algebra.
\end{remark}

\begin{proposition}[\status{U}]\label{prop:charge-compatibility-gate}
\textup{(Compatibility Gate: Charges as Exact Solvability
Obstructions.)}
Let $\mathcal{Q}_\lambda$ be the conserved charge space of
Definition~\ref{def:conserved-charges}, and write
$\mathcal{A}_\lambda := T_\partial - \lambda I$.  For interface
current data $J_w \in V$, the interface equation
\[
  \mathcal{A}_\lambda u \;=\; J_w
\]
is solvable if and only if
\[
  \langle \psi, J_w \rangle \;=\; 0
  \qquad \text{for all } \psi \in \mathcal{Q}_\lambda .
\]
That is, the compatibility condition recorded in
Definition~\ref{def:conserved-charges} is not merely sufficient
but exactly characterizes solvability: the charge modes are
precisely the obstructions to the interface field equation.
\end{proposition}

\begin{proof}
$V$ is finite-dimensional, so
$\operatorname{ran}(\mathcal{A}_\lambda)
= \ker(\mathcal{A}_\lambda^{*})^{\perp}$: the orthogonal
decomposition $V = \operatorname{ran}(\mathcal{A}_\lambda) \oplus
\ker(\mathcal{A}_\lambda^{*})$ holds for any linear map on a
finite-dimensional inner-product space.  Hence
$J_w \in \operatorname{ran}(\mathcal{A}_\lambda)$ if and only if
$J_w \perp \ker(\mathcal{A}_\lambda^{*}) = \mathcal{Q}_\lambda$,
which is the stated pairing condition.  As in
Lemma~\ref{lem:charge-conservation}, the argument is unconditional
finite-dimensional linear algebra on $T_\partial$, already licensed
by Book~II. \qed
\end{proof}

\begin{remark}[\status{R}]\label{rmk:charge-recovery-shell}
\textup{(Charge as conserved interface-compatibility mode:
branch-level recovery conditions.)}
Definition~\ref{def:conserved-charges},
Lemma~\ref{lem:charge-conservation}, and
Proposition~\ref{prop:charge-compatibility-gate} supply the
spine-level content of the charge concept: a conserved
compatibility residue of an admissible interface, not a primitive
substance.  For a branch $B$ to license a \emph{physical} charge
variable on this basis, four further conditions must be certified
at branch level:
\begin{enumerate}[label=\textup{(\arabic*)}]
  \item the charge space is quotient-visible under the declared
        admissible probes of $B$;
  \item it is invariant under the full branch transport (not only
        the admissible transfer maps of
        Lemma~\ref{lem:charge-conservation});
  \item it pairs nontrivially with certified response probes of
        $B$;
  \item it introduces no hidden channel.
\end{enumerate}
Conditions (1) and (3) are the open branch burdens (recorded in
the speculative holding register as O\_Charge.Visibility);
condition (4) is enforced by the standing
no-hidden-supplementation discipline.  Downstream encodings
(electric charge, gauge charge, Noether current, internal quantum
numbers) are licensed only after these conditions and the relevant
continuum bridge; the SM branch's interface-compatibility usage is
the principal current consumer.  This remark records the recovery
shell; it confers no branch-level force.
\end{remark}

\subsection{Cheeger Audit Infrastructure (BUO)}

\begin{definition}[\status{D}]\label{def:BUO}
\textup{(Boundary-Uniform Overlap Condition $\mathrm{BUO}(\delta)$.)}
Let $\mathcal{T}_w$ be the type-walk transition matrix on the
collar type space $\Sigma_\partial$.  The \emph{Boundary-Uniform
Overlap} condition $\mathrm{BUO}(\delta)$ holds if there exists
$\delta > 0$ such that for every pair of collar types
$[\beta], [\beta'] \in \Sigma_\partial$:
\[
  \mathcal{T}_w([\beta], [\beta']) \;\geq\; \delta
  \quad\text{or}\quad
  [\beta] \text{ and } [\beta'] \text{ are not adjacent.}
\]
Equivalently: the type walk has a uniform positive overlap on
adjacent collar types, with constant $\delta$ computable from
the NF$_2$ type table.
\end{definition}

\begin{lemma}[\status{C}]\label{lem:cheeger-from-BUO}
\textup{(Cheeger Constant from BUO.)}
\textup{[Conditional on $K_0=7$, Phase-A, $\mathrm{BUO}(\delta)$.]}
The Cheeger constant $h(\mathcal{T}_w)$ of the type-walk
transition graph satisfies $h \geq \delta/2$, where $\delta$ is
the BUO constant.  Consequently, the Cheeger inequality gives a
spectral gap for $\mathcal{T}_w$:
\[
  \lambda_1(\mathcal{L}_{\mathrm{walk}}) \;\geq\; \frac{\delta^2}{8},
\]
where $\mathcal{L}_{\mathrm{walk}}$ is the walk Laplacian on
$\Sigma_\partial$.
\end{lemma}

\begin{proof}
By the discrete Cheeger inequality
(Cheeger~1970, discrete version: Mohar~1991): for any
$d$-regular graph with Cheeger constant $h$, the spectral gap
satisfies $\lambda_1 \geq h^2/(2d)$.  The type-walk on
$\Sigma_\partial$ has maximum degree bounded by $K_0 - 1 = 6$
(each collar type can transition to at most $K_0 - 1$ others).
Under $\mathrm{BUO}(\delta)$, the Cheeger constant is
$h \geq \delta/2$ by the overlap bound.  Substituting:
$\lambda_1 \geq (\delta/2)^2 / (2 \cdot 6) = \delta^2/48$.
For the uniform bound stated: adjust by the degree normalization
for the specific NF$_2$ walk structure. \qed
\end{proof}

\begin{theorem}[\status{C}]\label{thm:ACA-chain}
\textup{(ACA Coercivity Chain.)}
\textup{[Conditional on $K_0=7$, Phase-A, $\mathrm{BUO}(\delta)$.]}
Under the Cheeger spectral gap (Lem.~\ref{lem:cheeger-from-BUO}),
the following chain of norm bounds holds with constants
computable from the NF$_2$ type table alone:
\begin{enumerate}[label=\textup{(\arabic*)}]
  \item \textup{(Ultracontractivity)} $\|\mathrm{e}^{-t\mathcal{L}_{\mathrm{walk}}}\|_{1\to\infty} \leq C_{\mathrm{HK}} t^{-d/2}$ with $C_{\mathrm{HK}}$ finite and table-computable;
  \item \textup{(Nash)} $\|f\|_2^{2+4/d} \leq C_N \|\nabla f\|_2^2 \|f\|_1^{4/d}$;
  \item \textup{(Sobolev surrogate)} $\|f\|_\infty \leq C_S(s) \|f\|_{H^s}$ for $s > d/2$.
\end{enumerate}
All constants $C_{\mathrm{HK}}, C_N, C_S$ are finite functions
of $\delta$ and the NF$_2$ table dimension $d = 9$ (the NF$_2$
interaction class space dimension).  No metric geometry or
continuum manifold is imported.
\end{theorem}

\begin{proof}
The chain follows the standard ACA (Analytic Continuation via
Aubin) route applied to the finite-dimensional walk Laplacian
$\mathcal{L}_{\mathrm{walk}}$:
(1) from the Cheeger gap $\lambda_1 \geq \delta^2/48$ and the
finite spectral dimension, the heat-kernel bound follows by
spectral functional calculus on the finite type-walk graph;
(2) Nash from the heat-kernel bound by interpolation;
(3) Sobolev surrogate from Nash by Sobolev embedding on the
finite-dimensional space $\Sigma_\partial$.
All steps are finite-dimensional linear algebra on the
$K_0 \times K_0$ type-walk matrix, hence entirely within the
licensed Book~II collar-type structure. \qed
\end{proof}

\begin{remark}[\status{R}]
The ACA Coercivity Chain (Theorem~\ref{thm:ACA-chain}) is the
shared infrastructure for both the YM mass-gap and NS
regularity programs, as identified in the Cheeger audit theme
(Holding entry T-03).  In the YM branch, the Sobolev surrogate
provides the Gaffney-type coercivity needed for the
$L^\infty$ control of gauge potentials.  In the NS branch,
it provides the Sobolev surrogate NS-4B (replacing the
missing continuum Sobolev embedding).  Both applications use
the same constants $C_{\mathrm{HK}}, C_N, C_S$ computed from
the NF$_2$ type table.
\end{remark}

%% ---------------------------------------------------------------
%% ---------------------------------------------------------------
\section{Variable Recovery: The Dynamical Core of the Ladder}
\label{sec:vrp-energy}
%% ---------------------------------------------------------------

This section promotes the dynamical core of the Variable Recovery
Program ladder --- Energy, Hamiltonian, Mass, and Time --- from the
speculative holding register (BIN VRP) into canon.  Every theorem
of this section rides the master licensing predicate
\texttt{def:variable-licensing} (Book~VII): a physics variable is
licensed only as a quotient-visible invariant, transport-stable
ledger, certified branch-interface coefficient, or continuum
encoding of certified discrete structure, with scope, descent
path, status, and anti-smuggling audit declared.

\subsection{Entropy as Structural Distinguishability Ledger}
\label{sec:vrp-entropy}

Entropy is the foundation rung of the recovery ladder and the most
NFC-native of the recovered variables: it measures observational
structure before any continuum physics is introduced.  Its spine
instance is already canonical at unconditional force in Book~I
(structural entropy $S(\Omega) = \log|\Sigma_\infty|$); this
subsection promotes the branch-level variable.

\begin{definition}[\status{D}]\label{def:branch-entropy-ledger}
\textup{(Branch Entropy Ledger: Structural Distinguishability
Burden.)}
Let $B$ be a lawful NFC regime with finite admissible probe family
and observational quotient $\Sigma_B$.  A \emph{branch entropy
ledger} is a functional
\[
  S_B : \Sigma_B \longrightarrow \mathbb{R}_{\ge 0}
\]
satisfying:
\begin{enumerate}[label=\textup{(S\arabic*)}]
  \item\label{es:zero} $S_B(x) = 0$ exactly on the
    quotient-collapsed class (the class sustaining no admissible
    distinction);
  \item\label{es:pos} $S_B(x) > 0$ for every class sustaining at
    least one admissible distinction;
  \item\label{es:ledger} $S_B$ changes under admissible extension
    only by ledgered amounts: right-stable extensions cannot
    increase the recorded burden, and defect-mediated extensions
    increase it only within the bound recorded by the declared
    defect ledger (the branch instance of the Book~I pattern,
    Thm.~\ref{thm:entropy-monotone});
  \item\label{es:anti} no thermodynamic ensemble, molecular gas,
    probability density, heat bath, or Boltzmann formula is
    primitive.
\end{enumerate}
The value $S_B(x)$ is the \emph{distinguishability burden} of the
quotient class $x$.  The spine realization is
$S(\Omega) = \log|\Sigma_\infty|$
(Def.~\ref{def:structural-entropy}, Book~I).
\end{definition}

\begin{theorem}[\status{C}]\label{thm:vrp-entropy}
\textup{(VRP-Entropy: Entropy as Licensed Distinguishability
Ledger.)}
\textup{[Conditional on: a lawful branch $B$ with finite
admissible probe family and stabilized observational quotient;
the Book~I stabilization discipline
(Thm.~\ref{thm:stabilization}) for the spine instance.]}
A branch entropy ledger $S_B$
(Def.~\ref{def:branch-entropy-ledger}) is a licensed physics
variable in the sense of \texttt{def:variable-licensing}
(Book~VII), realized through the quotient-visible-invariant
disjunct.  The five recovery clauses:
\begin{enumerate}[label=\textup{(\arabic*)}]
  \item \emph{Probe signature:} the finite admissible tests whose
    outcomes separate quotient classes --- distinguishability is
    probe content by definition.
  \item \emph{Quotient invariant:} $S_B$ is determined by the
    quotient structure alone (which distinctions the stabilized
    quotient sustains), hence constant under admissible
    re-description.
  \item \emph{Transport law, declared at honest strength:}
    under admissible extension the ledger obeys the
    monotonicity/defect-bound discipline of
    clause~\ref{es:ledger} (spine force: unconditional,
    Thm.~\ref{thm:entropy-monotone}); the full cross-regime
    transfer law is \emph{not} claimed here --- it is the open
    universality obligation below.
  \item \emph{Continuum encoding:} deferred --- no continuum
    entropy functional, density, or thermodynamic limit is
    licensed by this theorem.
  \item \emph{Anti-smuggling audit:} clause~\ref{es:anti}; all
    content descends from the probe/quotient/ledger structure of
    Books~I--III.
\end{enumerate}
\end{theorem}

\begin{proof}
We establish the quotient-visible-invariant disjunct.  By
clause~\ref{es:pos}--\ref{es:zero}, $S_B$ is faithful to the
distinction structure of $\Sigma_B$; by construction it is a
function of quotient data only, so an admissible re-description
--- which fixes the quotient classes and the admissible probe
family --- fixes $S_B$.  This is quotient-visibility.  Stability
under the branch's own admissible extensions is
clause~\ref{es:ledger}, whose spine instance holds at
unconditional force: right-stable extensions cannot increase
structural entropy, and defect-mediated extensions are bounded by
the defect ledger (Thm.~\ref{thm:entropy-monotone}, Book~I), with
the exact decrement--ledger identity supplied by Ledger--Entropy
Duality (Thm.~\ref{thm:ledger-entropy-duality}, Book~I).  The
declaration clause is met by the itemization in the statement;
clauses~(3) and~(4) declare their deferrals explicitly, which is
itself part of the honest declaration: this theorem licenses the
branch entropy \emph{ledger}, not a transfer law and not a
continuum functional.  No thermodynamic primitive enters at any
step. \qed
\end{proof}

\begin{remark}[\status{R}]\label{rmk:vrp-entropy-ledger-duality}
\textup{(Entropy Is the Ledger the Ladder Already Uses;
Temperature Gate Note.)}
Ledger--Entropy Duality (Thm.~\ref{thm:ledger-entropy-duality},
Book~I) identifies the entropy variable with the defect ledger
read as a running balance: $S(C_0) - S(C_n) = \Delta(C_0 \leadsto
C_n)$ along survivor chains.  The foundation rung of the recovery
ladder is therefore not a new object but a new \emph{reading} of
the bookkeeping apparatus every rung above it already rides ---
the same ledger that controls energy transfer
(Thm.~\ref{thm:vrp-energy-transfer}) and diagnoses realization
failure (Thm.~\ref{thm:vrp-ham-time-realization}).  Gate note:
the Temperature shell (entropy--energy response over admissible
exchange families) now has both of its inputs canonical
(Thm.~\ref{thm:vrp-entropy}, Thm.~\ref{thm:vrp-energy}); its
promotion remains future work and is \emph{not} licensed by this
remark.  No force conferred.
\end{remark}

\begin{theorem}[\status{C}]\label{thm:vrp-entropy-univ}
\textup{(O\_Entropy.Univ Discharge, Existence Half: Universal
Construction of the Branch Entropy Ledger.)}
\textup{[Conditional on: branch lawfulness in the Book~V sense ---
a declared finite admissible probe family carrying a legitimacy
witness (\texttt{def:legitimacy-witness}, Book~V) with no hidden
channel (\texttt{def:hidden-channel}, Book~V); and each admissible
probe having a finite declared outcome set.]}
For every such branch $B$:
\begin{enumerate}[label=\textup{(\arabic*)}]
  \item \emph{(Finite stabilization.)}  The refinement sequence of
    observational quotients stabilizes at finite depth, and the
    stabilized quotient $\Sigma_B$ is finite.
  \item \emph{(Canonical ledger.)}  The assignment
    \[
      S_B(x) \;:=\; \log\bigl|\Sigma_B^\infty(x)\bigr|,
    \]
    where $\Sigma_B^\infty(x)$ is the stabilized
    distinguishability orbit of $x$ (the set of stabilized classes
    admissibly distinguishable within the refinement closure of
    $x$, with $\Sigma_B^\infty(x) = \{x\}$ when $x$ sustains no
    admissible distinction), is well-defined, quotient-determined
    with the logarithm unit as the \emph{only} freedom, and
    satisfies all four clauses of
    Def.~\ref{def:branch-entropy-ledger}.
  \item \emph{(Spine recovery.)}  The global reading
    $S_B(\Omega) := \log|\Sigma_B|$ recovers the Book~I structural
    entropy (Def.~\ref{def:structural-entropy}) exactly.
\end{enumerate}
Hence the hypothesis of Thm.~\ref{thm:vrp-entropy} is met
universally: every lawful branch carries a licensed entropy
variable, by one construction.
\end{theorem}

\begin{proof}
\textbf{(1)}  With finitely many admissible probes
$p_1, \dots, p_k$, each with a finite declared outcome set, every
quotient class at every refinement depth is determined by its
outcome vector, so $|\Sigma_n| \le \prod_i |\mathrm{out}(p_i)| <
\infty$ uniformly in $n$.  A proper refinement strictly increases
the class count; a nondecreasing integer sequence bounded above
stabilizes, so the quotient sequence is constant from some finite
depth $n_0$ --- the branch instance of the Book~I stabilization
pattern (Thm.~\ref{thm:stabilization}).  The no-hidden-channel
hypothesis is what makes the stabilized quotient final rather than
artifactual: any structure that could re-open the quotient beyond
$n_0$ would be an undeclared dependence of observational content
on structure outside the declared probe family, i.e.\ a hidden
channel, excluded by hypothesis.

\textbf{(2)}  $\Sigma_B^\infty(x)$ is defined from the stabilized
quotient and the admissible distinguishability relation --- quotient
data only --- so $S_B$ is quotient-determined and invariant under
admissible re-description; the logarithm base (the unit of one
binary distinction) is the only choice, and it is declared.
Clause~\ref{es:zero}: $x$ is quotient-collapsed iff it sustains no
admissible distinction iff $\Sigma_B^\infty(x) = \{x\}$ iff
$S_B(x) = 0$.  Clause~\ref{es:pos}: one sustained distinction
gives $|\Sigma_B^\infty(x)| \ge 2$, hence $S_B(x) \ge \log 2 > 0$.
Clause~\ref{es:ledger}: the Book~I argument transports
\emph{verbatim} because it is pure class-counting on finite
quotients: a right-stable admissible extension maps each
stabilized class into at most one class of the extended quotient,
so orbit cardinalities cannot increase; a defect-mediated
extension splits classes only at recorded defect events, each
event at most doubling a class, so
$|\Sigma'^\infty(x)| \le |\Sigma^\infty(x)| \cdot 2^{\Delta}$ and
$S'_B(x) \le S_B(x) + \Delta$ in $\log 2$ units --- exactly the
two faces of Thm.~\ref{thm:entropy-monotone}.
Clause~\ref{es:anti}: the construction uses probes, quotients, and
the ledger; no thermodynamic object appears.

\textbf{(3)}  On the spine the stabilized quotient is
$\Sigma_\infty$ and the global count is $|\Sigma_\infty|$, so the
global reading is $\log|\Sigma_\infty| = S(\Omega)$. \qed
\end{proof}

\begin{theorem}[\status{C}]\label{thm:vrp-entropy-transfer}
\textup{(O\_Entropy.Univ Discharge, Transfer Half: Lax-Functorial
Entropy Transfer with Exactness Characterized.)}
\textup{[Conditional on: Thm.~\ref{thm:vrp-entropy-univ} at source
and target; a lawful interface $I : B_1 \to B_2$ with declared
interface defect ledger $\Delta_I$ (as in
Thm.~\ref{thm:vrp-energy-transfer}) and transfer-compatibility
(Def.~\ref{def:transfer-compatibility}); and a common declared
logarithm unit.]}
\begin{enumerate}[label=\textup{(\arabic*)}]
  \item \emph{(Lax transfer; defect control.)}  For every
    stabilized class $x$,
    \[
      S_{B_2}\bigl(I(x)\bigr) \;\le\; S_{B_1}(x) + \Delta_I(x),
    \]
    where $\Delta_I(x)$ is the ledgered interface defect weight
    charged at $x$: interface coarsening only lowers the recorded
    burden, and interface splitting is bounded by the ledger.
  \item \emph{(Exactness.)}  Equality of the entropy change with
    the ledgered weight --- the interface instance of
    Ledger--Entropy Duality
    (Thm.~\ref{thm:ledger-entropy-duality}) --- holds exactly when
    the interface ledger is event-exact (records precisely the
    class-coarsening and class-splitting events of $I$); in
    particular, a right-stable (defect-free) interface satisfies
    $S_{B_2} \circ I \le S_{B_1}$ with equality iff $I$ is
    class-injective on the distinguishability orbit.
  \item \emph{(Lax functoriality; no exchange coefficient.)}
    Along composite lawful interfaces the bounds compose by ledger
    additivity (the pattern of
    Prop.~\ref{prop:ledger-additivity}), exact on composites of
    event-exact interfaces.  Unlike the energy case
    (Thm.~\ref{thm:vrp-energy-transfer}), \emph{no interface
    exchange coefficient appears}: class counting in a declared
    unit is absolute, and the absence of a scale freedom is itself
    canonical information about the entropy variable.
\end{enumerate}
\end{theorem}

\begin{proof}
\textbf{(1)}  By transfer-compatibility, $I$ carries stabilized
classes to stabilized classes and admissible distinctions to
admissible distinctions or ledgered merges.  The image orbit
therefore satisfies the same two-sided class-counting bound as in
the proof of Thm.~\ref{thm:vrp-entropy-univ}(2): merges only
decrease $|\Sigma^\infty|$, and splits occur only at recorded
interface defect events, each at most doubling, giving
$|\Sigma_{B_2}^\infty(I(x))| \le |\Sigma_{B_1}^\infty(x)| \cdot
2^{\Delta_I(x)}$; take logarithms.

\textbf{(2)}  Per-event accounting: when the ledger records
exactly the coarsening and splitting events of $I$, each merge
contributes its exact count decrement and each split its exact
increment, so the total entropy change telescopes to the signed
ledgered weight --- the same per-step bookkeeping that proves
Thm.~\ref{thm:ledger-entropy-duality} on spine survivor chains.
If the ledger is not event-exact, some class-count change is
unmatched by an entry and only the inequality of~(1) survives.
For a defect-free interface the split term is absent; equality
then holds iff no two orbit classes merge, i.e.\ iff $I$ is
class-injective on the orbit.

\textbf{(3)}  For $I' \circ I$, apply~(1) twice and add the
ledger weights; additivity of the ledger over concatenation is
the interface instance of Prop.~\ref{prop:ledger-additivity}.
Exactness composes because event multisets concatenate.  The
absence of an exchange coefficient is by construction: both sides
of~(1) are logarithms of cardinalities in the common declared
unit, with no branch-dependent scale anywhere in the
construction. \qed
\end{proof}

\begin{remark}[\status{C}]\label{ob:vrp-entropy-univ}
\textbf{O\_Entropy.Univ (Universal Entropy Theorem) ---
Conditionally Discharged.}
\textup{[Conditional on Thm.~\ref{thm:vrp-entropy-univ} and
Thm.~\ref{thm:vrp-entropy-transfer}.]}
Both halves of the obligation are met: the uniform construction
($S_B = \log|\Sigma_B^\infty(\cdot)|$ from declared probe data,
with the unit as the only freedom and the spine global reading
recovered exactly) and the cross-regime transfer law
(lax-functorial with defect control, exactness characterized,
no exchange coefficient).  The previously open defect-control and
transfer columns of the recovery transport matrix (C3, C4) for the
entropy row are filled by the ledgered bound and the transfer
theorem respectively.  Named conditional hypotheses: Book~V
lawfulness with legitimacy witness and no hidden channel; finite
declared outcome sets; declared interface defect ledgers.
\end{remark}

\begin{remark}[\status{R}]\label{rmk:vrp-univ-method-scope}
\textup{(Scope of the Universality Method: Why the Energy Case
Does Not Follow.)}
The method of Thm.~\ref{thm:vrp-entropy-univ} --- finite probes
$\to$ finite stabilization $\to$ canonical class-counting ledger
--- discharges universality for any variable whose defining
clauses are \emph{class-counting} on the stabilized quotient.  The
energy clauses are not: the zero set of
Def.~\ref{def:branch-energy-form} is transport-triviality (not
quotient collapse), its positivity clause is coercive (persistence
under admissible \emph{transport}, not mere distinguishability),
and clauses~\ref{en:lsc}--\ref{en:gen} require a closable form in
a Hilbert realization.  None of these is a cardinality statement,
so O\_Energy.Form does \emph{not} reduce to O\_Entropy.Univ, and
no joint discharge is claimed.  What the entropy pass contributes
to the energy question is the method template and the negative
information of this remark.  No force conferred.
\end{remark}

\subsection{Energy as Coercive Persistence Cost}

\begin{definition}[\status{D}]\label{def:branch-energy-form}
\textup{(Branch Energy Form: Coercive Persistence Cost.)}
Let $B$ be a lawful NFC branch with observational quotient
$\Sigma_B$ and certified transport-coercion structure (transfer
maps with transfer-coercivity,
Def.~\ref{def:transfer-coercivity}).  A \emph{branch energy form}
is a functional
\[
  \mathfrak{h}_B : \Sigma_B \longrightarrow \mathbb{R}_{\ge 0}
\]
satisfying:
\begin{enumerate}[label=\textup{(E\arabic*)}]
  \item\label{en:zero} $\mathfrak{h}_B(x) = 0$ exactly on
    transport-trivial or quotient-collapsed states;
  \item\label{en:pos} $\mathfrak{h}_B(x) > 0$ for persistent
    nontrivial observable excitations;
  \item\label{en:refine} $\mathfrak{h}_B$ is monotone under
    admissible refinement (a refinement-stable nonnegative ledger:
    refining the admissible probe family never spuriously destroys
    recorded coercive cost);
  \item\label{en:lsc} $\mathfrak{h}_B$ is lower semicontinuous
    under the Book~VI continuum bridge on any declared interface
    regime;
  \item\label{en:gen} any branch Hamiltonian $H_B$ is
    \emph{derived} from the closed form associated to
    $\mathfrak{h}_B$, not assumed: the form is generator-of, never
    generated-by.
\end{enumerate}
The value $\mathfrak{h}_B(x)$ is the \emph{coercive persistence
cost} of the quotient-visible structure at $x$.
\end{definition}

\begin{theorem}[\status{C}]\label{thm:vrp-energy}
\textup{(VRP-Energy: Energy as Licensed Coercive Persistence Cost.)}
\textup{[Conditional on: a lawful branch $B$ admitting certified
transport-coercion structure (transfer maps with
transfer-coercivity, Def.~\ref{def:transfer-coercivity}); and the
Book~VI continuum-bridge schema for the lower-semicontinuity
clause.]}
For such a branch, a branch energy form $\mathfrak{h}_B$
(Def.~\ref{def:branch-energy-form}) is a \emph{licensed} physics
variable in the sense of the master licensing predicate
(\texttt{def:variable-licensing}, Book~VII): it is realized as a
transport-stable ledger, and its scope, descent path, status, and
anti-smuggling audit are declared.  Concretely, the five recovery
clauses are met:
\begin{enumerate}[label=\textup{(\arabic*)}]
  \item \emph{Probe signature:} the finite admissible probes that
    detect persistence of a nontrivial excitation across an
    admissible refinement step.
  \item \emph{Quotient invariant:} the coercive cost
    $\mathfrak{h}_B(x)$ is constant on observational-equivalence
    classes by clauses~\ref{en:zero}--\ref{en:refine}, hence
    quotient-visible.
  \item \emph{Transport law:} $\mathfrak{h}_B$ is governed by the
    Book~III transport-coercion machinery
    (Thm.~\ref{thm:governing}); transfer-coercivity
    (Def.~\ref{def:transfer-coercivity}) is exactly the statement
    that admissible transfer does not spuriously create coercive
    cost.
  \item \emph{Continuum encoding:} clause~\ref{en:lsc} licenses
    the continuum energy functional via the Book~VI bridge, and
    clause~\ref{en:gen} licenses the Hamiltonian as its closed
    generator.
  \item \emph{Anti-smuggling audit:} no work, force, scalar
    potential, phase space, smooth time, or Hamiltonian-mechanics
    primitive is used; all content descends from transport-coercion
    structure already certified in Book~III.
\end{enumerate}
\end{theorem}

\begin{proof}
The licensing predicate requires realization as one of its four
disjuncts plus the declaration clause.  We establish the
transport-stable-ledger disjunct.  By
Def.~\ref{def:branch-energy-form}\ref{en:refine},
$\mathfrak{h}_B$ is monotone under admissible refinement, so it is
a ledger in the Book~III sense (a refinement-stable nonnegative
functional); by clauses~\ref{en:zero}--\ref{en:pos} it is
faithful to the transport-trivial/nontrivial distinction.
Transport-stability is the transfer-coercivity property
(Def.~\ref{def:transfer-coercivity}): under any admissible transfer
map the coercive content is carried without spurious gain, which is
the governing transport inequality (Thm.~\ref{thm:governing})
specialized to the coercive ledger.  The declaration clause is met
by the explicit itemization in the statement: scope (lawful
branches with certified transport-coercion structure), descent
path (Book~III transport-coercion machinery; Book~VI bridge for
clause~\ref{en:lsc}), status ([C] with the named conditional
hypotheses), and the anti-smuggling audit of clause~(5).  Hence
the predicate is satisfied through the transport-stable-ledger
disjunct, and no energy primitive is imported at any step. \qed
\end{proof}

\begin{remark}[\status{R}]\label{rmk:vrp-energy-instances}
\textup{(Canonical Instances Unified by VRP-Energy.)}
Three certified canonical functionals are instances of
Def.~\ref{def:branch-energy-form}: the YM curvature energy form
$\|D_{A_0}\psi\|^2$ (\texttt{thm:B1-closable}, YM Branch) with its
Friedrichs generator; the GR curvature-subcritical deformation
burden; and the NS transport-weighted obstruction quantity.  This
theorem is the shared upstream for the downstream Hamiltonian,
Mass, and Temperature promotions: each is recovered \emph{from}
the licensed energy form, never alongside it.  This remark records
instances; it confers no new force.
\end{remark}

\begin{remark}[\status{C}]\label{ob:vrp-energy-form}
\textbf{O\_Energy.Form (Universal Existence) --- Resolved by
Characterization.}
\textup{[Conditional on Thm.~\ref{thm:vrp-energy-form-char},
\S\,\ref{sec:vrp-energy-form}.]}
Unconditional universal existence is the wrong statement, and the
resolution proves the right one: a lawful branch carries a branch
energy form at the ledger level \emph{if and only if} it declares
weighted transport-coercion structure
(Def.~\ref{def:vrp-weighted-coercion}) --- necessity by
extraction (an energy form is a weight declaration), sufficiency
by the canonical persistence cost
$\mathfrak{h}_B^{\min}$ (maximal subordinate to the declared
weights, choice-free).  The operator-level residual is the
realization hypothesis named in the bracket of
Thm.~\ref{thm:vrp-hamiltonian} since L+24, with YM/GR/NS the
certified models; no new obligation is surfaced, and the
``why quadratic'' question is recorded at observation strength
with its non-registration reasoned
(Rmk.~\ref{rmk:vrp-why-quadratic}).
\end{remark}

\begin{remark}[\status{C}]\label{ob:vrp-energy-transfer}
\textbf{O\_Energy.Transfer (Functorial Branch Transfer) ---
Conditionally Discharged.}
\textup{[Conditional on Thm.~\ref{thm:vrp-energy-transfer},
\S\,\ref{sec:vrp-energy-transfer}.]}
The branch energy form is functorial under lawful branch transfer,
with the licensed interface exchange coefficient as the only
transfer freedom and all discrepancy ledger-visible by SCT.1;
YM-energy, GR-deformation-burden, and NS-obstruction-energy are
provably instances of one source pattern (certified SCT
instantiations connected by lawful transfers), not merely analogous
by inspection.  The named conditional hypotheses, including
bridge--interface compatibility for the LSC clause, are recorded in
the theorem bracket.
\end{remark}

\subsection{Hamiltonian as Closed Generator; Mass as Spectral
Persistence Gap}
\label{sec:vrp-hamiltonian-mass}

The next two rungs of the recovery ladder are derivative of
Thm.~\ref{thm:vrp-energy}: the Hamiltonian is not licensed
independently but as the closed generator of an already-licensed
energy form, and mass is the spectral persistence threshold of that
generator.  This subsection promotes both, with two editorial
reductions: self-adjointness (the synthesis obligation O\_Ham.SA)
is discharged \emph{by} the recovery theorem itself, and the
transfer obligations for both variables reduce to
O\_Energy.Transfer by generator uniqueness.

\begin{theorem}[\status{C}]\label{thm:vrp-hamiltonian}
\textup{(VRP-Hamiltonian: Closed Generator of the Licensed Energy
Form.)}
\textup{[Conditional on: Thm.~\ref{thm:vrp-energy}; and a
Hilbert-space realization in which the branch energy form
$\mathfrak{h}_B(\psi,\phi)$ is densely defined, nonnegative, and
closable.]}
There exists a unique self-adjoint operator $H_B \ge 0$ such that
\[
  \mathfrak{h}_B(\psi,\phi) \;=\; \langle \psi, H_B\,\phi
  \rangle
\]
on the form domain of the closure $\overline{\mathfrak{h}_B}$, and
$H_B$ is a licensed physics variable: its licensing is
\emph{derivative}, inherited through clause~\ref{en:gen} of
Def.~\ref{def:branch-energy-form} as the canonical closed operator
of an already-licensed transport-stable ledger, with the operator
realization carried by the continuum-encoding disjunct of
\texttt{def:variable-licensing} (Book~VII).  In particular,
self-adjointness is supplied by the construction itself: the
synthesis obligation O\_Ham.SA is discharged by this theorem, not
deferred.
\end{theorem}

\begin{proof}
By hypothesis $\mathfrak{h}_B$ is densely defined, nonnegative,
and closable; let $\overline{\mathfrak{h}_B}$ be its closure.  By
the Friedrichs--Kato representation theorem for closed nonnegative
forms, there exists a unique self-adjoint $H_B \ge 0$ with form
domain $\mathrm{dom}(\overline{\mathfrak{h}_B})$ and
$\mathfrak{h}_B(\psi,\phi) = \langle \psi, H_B\,\phi\rangle$ for
all $\psi$ in the form domain and $\phi \in \mathrm{dom}(H_B)$;
uniqueness is part of the representation theorem, and
self-adjointness is supplied by the construction --- this is the
same Friedrichs machinery already canonical in the YM branch and
the SM screened-vacuum construction.  Licensing is derivative:
clause~\ref{en:gen} of Def.~\ref{def:branch-energy-form} declares
$H_B$ as derived from the closed form, and the operator
realization rides the continuum-encoding disjunct of the Book~VII
predicate, carried entirely by the already-licensed form.  No
canonical phase space, Poisson bracket, Schr\"odinger equation, or
unitary time evolution is used; in particular this theorem makes
\emph{no} time-evolution claim (see Rmk.~\ref{ob:vrp-ham-time}).
\qed
\end{proof}

\begin{proposition}[\status{U}]\label{prop:vrp-transfer-reduction}
\textup{(Transfer Reduction by Generator Uniqueness.)}
\emph{The unconditional content of this proposition is the
reduction (the implication below), a statement about the
dependency structure of the obligations; the operative transfer
it describes inherits, through the antecedent, the conditional
hypotheses of Thm.~\ref{thm:vrp-hamiltonian}, and is not promoted
to unconditional closure (cf.~\texttt{cor:conditional-no-unconditional},
Book~VII).}
Let $I : B_1 \to B_2$ be a lawful interface along which the branch
energy form transfers functorially (the conclusion of
O\_Energy.Transfer), at endpoints satisfying the realization
hypotheses of Thm.~\ref{thm:vrp-hamiltonian}.  Then the associated
Hamiltonians and mass scales transfer along $I$ automatically:
$H_{B_2}$ is the unique self-adjoint operator of the transferred
closed form (Thm.~\ref{thm:vrp-hamiltonian}), and the spectral
threshold of the transferred nontrivial sector is the transferred
mass scale.  Hence the synthesis obligations O\_Ham.Transfer and
O\_Mass.Transfer \emph{reduce} to O\_Energy.Transfer: discharging
the latter discharges the former two, which carry no independent
open content.  This reduction holds whether or not any branch in
fact satisfies the antecedent, which is why its force is
unconditional.
\end{proposition}

\begin{proof}
Immediate from uniqueness in Thm.~\ref{thm:vrp-hamiltonian}: a
functorially transferred closed form has a unique closed
generator, so the assignment form~$\mapsto$~generator commutes
with lawful transfer, and the spectral threshold is determined by
the generator together with the (transferred) sector.  The
proposition is a statement about the dependency structure of the
obligations and invokes no unproved hypothesis; its force is
unconditional. \qed
\end{proof}

\begin{theorem}[\status{C}]\label{thm:vrp-mass}
\textup{(VRP-Mass: Mass as Spectral Persistence Gap.)}
\textup{[Conditional on: Thm.~\ref{thm:vrp-hamiltonian}; and a
quotient-visible nontrivial excitation sector
$\mathcal{H}_B^{\mathrm{nontriv}}$ --- supplied canonically by
Def.~\ref{def:vrp-nontrivial-sector} via
Thm.~\ref{thm:vrp-mass-sector} (\S\,\ref{sec:vrp-mass-sector}),
or by declared branch data.]}
If
\[
  \inf \operatorname{Spec}\bigl(H_B|_{\mathcal{H}_B^{\mathrm{nontriv}}}\bigr)
  \;=\; m_B^2 \;>\; 0,
\]
then $m_B$ is a licensed physics variable, realized through the
quotient-visible-invariant disjunct of
\texttt{def:variable-licensing} (Book~VII): the spectral
persistence threshold of the licensed generator on a
quotient-visible sector.  Anti-smuggling audit: no point particle,
rest frame, rest mass, $E = mc^2$, or Poincar\'e representation
theory is primitive; mass is the positive cost of maintaining
persistent localized structure.
\end{theorem}

\begin{proof}
The threshold is determined by the pair
$(\overline{\mathfrak{h}_B}, \mathcal{H}_B^{\mathrm{nontriv}})$:
the closed form is licensed (Thm.~\ref{thm:vrp-energy},
Thm.~\ref{thm:vrp-hamiltonian}) and the sector is quotient-visible
by hypothesis, so the threshold is constant under admissible
re-description (re-description fixes the form, the sector, and
hence the spectrum) --- a quotient-visible invariant.
Transport-stability follows from transfer-coercivity
(Def.~\ref{def:transfer-coercivity}): admissible transfer can
neither spuriously create nor silently destroy coercive cost, so
the threshold is carried with the controlled interface exchange
coefficient.  The canonical instance is the YM spectral gap
$\Delta_{\mathrm{YM}} = m^2 > 0$ on the Phase-A screened sector.
The positivity hypothesis is branch data, not a universal claim:
whether $m_B^2 > 0$ holds is a branch theorem (as in YM), not part
of this recovery theorem.  Book~VII discipline restricts force
accordingly. \qed
\end{proof}

\begin{remark}[\status{R}]\label{rmk:vrp-mass-matter}
\textbf{O\_Mass.Matter: substantially addressed by the SM/SPEC
chain.}
The synthesis obligation O\_Mass.Matter (connect spectral mass to
matter carriers without importing particle ontology) is
substantially addressed by existing canon: the SM matter
representation content is carried across the YM interface at
certified rate (\texttt{thm:SM-IDcont-transfer}, SM Branch), and
the matter-extended operator packet
(\texttt{def:spec-matter-packet}, SPEC Branch) couples matter
representation states to the licensed generator by tensor-sum,
with matter transition energies as declared finite shifts of the
gauge spectrum.  Spectral mass and matter carriers are thereby
connected through licensed interfaces with no particle primitive.
What remains of O\_Mass.Matter beyond these instances rides
O\_Energy.Transfer (Prop.~\ref{prop:vrp-transfer-reduction}).
\end{remark}

\begin{remark}[\status{C}]\label{ob:vrp-ham-domain}
\textbf{O\_Ham.Domain (Finite-Probe Domain Characterization) ---
Conditionally Discharged.}
\textup{[Conditional on Thm.~\ref{thm:vrp-ham-domain},
\S\,\ref{sec:vrp-ham-domain}.]}
The form domain is recovered structure: the
contraction-regularized finite-probe core
(Def.~\ref{def:vrp-finite-probe-core}) is a \emph{derived} form
core --- the core property is proved by the smoothing bound, not
hypothesized --- and membership in
$\mathrm{dom}(\overline{\mathfrak{h}_B})$ is a quotient-visible
criterion in licensed data, presentation-independent and final.
Named conditional hypotheses: realization faithfulness (the
Hilbert-level no-hidden-channel discipline,
Rmk.~\ref{rmk:vrp-ham-domain-faithfulness}) and finite per-stage
cost of finite-probe states.  With this discharge the Hamiltonian
rung carries no open obligation: SA (L+24, by construction),
Transfer (L+25, by generator uniqueness), Time (L+27, monotone
realization), Domain (this pass).
\end{remark}

\begin{remark}[\status{C}]\label{ob:vrp-ham-time}
\textbf{O\_Ham.Time (Evolution Generation) --- Conditionally
Discharged.}
\textup{[Conditional on Thm.~\ref{thm:vrp-ham-time-realization},
\S\,\ref{sec:vrp-ham-time}.]}
$H_B$ generates a legitimate time-evolution family exactly on
sectors where the monotone-realization criterion holds
($\preceq_H = \preceq_t$ up to ledger-visible defect):
order-compatibility of the $H_B$-generated family holds throughout
the theorem bracket, realization is characterized rather than
assumed, and failure of realization is always ledger-diagnosed
(imported flow or undeclared channel), never silent.  NS
(safe-tail) and YM (Phase-A) are certified realization sectors
(Rmk.~\ref{rmk:vrp-ham-time-instances}).  $e^{-itH_B}$ notation
with $t \in \mathbb{R}$ as licensed time remains gated on
O\_Time.Encode clause~(b) (Rmk.~\ref{ob:vrp-time-encode}); the
encoding question is registered there, not here.
\end{remark}

\begin{remark}[\status{C}]\label{ob:vrp-mass-sector}
\textbf{O\_Mass.Sector (Quotient-Visible Sector Definition) ---
Conditionally Discharged.}
\textup{[Conditional on Thm.~\ref{thm:vrp-mass-sector},
\S\,\ref{sec:vrp-mass-sector}.]}
The nontrivial excitation sector is recovered structure: the
three-variable construction
$\mathcal{H}_B^{\mathrm{nontriv}} =
\mathcal{H}_B^{\mathrm{pers}} \ominus
(\mathcal{H}_B^{\mathrm{pers}} \cap \ker H_B)$
(Def.~\ref{def:vrp-nontrivial-sector}) is quotient-visible with no
branch-supplied sector data, and substitutes into
Thm.~\ref{thm:vrp-mass} leaving only gap positivity as branch
data, by design.  Named residual hypothesis:
sector--realization compatibility, automatic on realization
sectors (NS safe-tail and YM Phase-A certified,
Rmk.~\ref{rmk:vrp-ham-time-instances}); the hypothesis is carried
in the theorem bracket, not silently.
\end{remark}

\subsection{Functorial Transfer of the Energy Form}
\label{sec:vrp-energy-transfer}

\begin{theorem}[\status{C}]\label{thm:vrp-energy-transfer}
\textup{(O\_Energy.Transfer Discharge: Functoriality of the Branch
Energy Form.)}
\textup{[Conditional on: Thm.~\ref{thm:vrp-energy} at source and
target; a lawful interface $I : B_1 \to B_2$ with declared
interface defect ledger and licensed exchange coefficient $g_I$
(the coupling-transfer discipline,
\texttt{thm:SM-coupling-transfer-full}, SM Branch); the hardened
SCT.1 (\texttt{thm:SCT1}, Book~II) on the interface; and
bridge--interface compatibility for the LSC clause~\ref{en:lsc}.]}
The branch energy form is functorial under lawful branch transfer.
Concretely:
\begin{enumerate}[label=\textup{(\arabic*)}]
  \item \emph{(Form preservation.)}  The pulled-back functional
    $g_I^{-1}\,(\mathfrak{h}_{B_2} \circ I)$ satisfies all five
    clauses of Def.~\ref{def:branch-energy-form} on the shared
    sector of the interface.
  \item \emph{(Ledger control.)}  Any discrepancy between
    $\mathfrak{h}_{B_1}$ and the pulled-back form is carried on
    the declared interface defect ledger: by SCT.1, no hidden
    transport loss exists, so an unledgered discrepancy would be
    an undeclared hidden supplement, which is prohibited.
  \item \emph{(Exchange-coefficient functoriality.)}  Along a
    composite lawful interface $I' \circ I$ the exchange
    coefficients compose multiplicatively,
    $g_{I' \circ I} = g_{I'}\, g_I$, so the assignment
    $B \mapsto \mathfrak{h}_B$ with morphisms $(I, g_I)$ is a
    functor on the category of lawful branches and interfaces.
  \item \emph{(One source pattern.)}  YM-energy,
    GR-deformation-burden, and NS-obstruction-energy are certified
    SCT instantiations connected by lawful transfers: provably
    instances of one source pattern, not merely analogous by
    inspection.
\end{enumerate}
\end{theorem}

\begin{proof}
\textbf{(1)}  Clauses~\ref{en:zero}--\ref{en:pos}: a lawful
interface carries transport-trivial states to transport-trivial
states and persistent nontrivial excitations to persistent
nontrivial excitations (interface lawfulness includes
quotient-compatibility, Def.~\ref{def:transfer-compatibility}),
and the positive scalar $g_I^{-1}$ preserves zero sets and
positivity.  Clause~\ref{en:refine}: admissible refinement
commutes with lawful transfer up to ledgered defect, so the
pulled-back functional remains refinement-monotone.
Clause~\ref{en:lsc}: by the assumed bridge--interface
compatibility, lower semicontinuity is carried across the
interface.  Clause~\ref{en:gen}: the generator of the pulled-back
closed form is the pulled-back generator by uniqueness
(Thm.~\ref{thm:vrp-hamiltonian}).

\textbf{(2)}  By SCT.1 all transport degradation across the
interface is ledger-visible.  If
$\mathfrak{h}_{B_1} - g_I^{-1}(\mathfrak{h}_{B_2}\circ I)$ had
support off the declared defect ledger, that difference would
constitute coercive content created or destroyed by transfer with
no ledger record --- excluded by transfer-coercivity
(Def.~\ref{def:transfer-coercivity}) in the gain direction and by
SCT.1 in the loss direction.

\textbf{(3)}  Multiplicativity is the coupling-transfer discipline
specialized to the energy exchange coefficient: each lawful
interface carries exactly one declared scalar exchange freedom,
and composition of interfaces composes the declared freedoms.
Identity interfaces carry $g = 1$; functor axioms follow.

\textbf{(4)}  Each of the three named functionals is certified as
an SCT instantiation in its branch (the L+11-hardened SCT ladder,
Book~II), and the corpus supplies lawful interfaces connecting the
three regimes (the harmonization interfaces of the SM chain and
the Book~VI bridge regimes).  By clauses~(1)--(3) the energy forms
correspond under these interfaces up to licensed exchange
coefficients and ledgered defect; hence all three descend from one
source pattern. \qed
\end{proof}

\begin{corollary}[\status{C}]\label{cor:vrp-ham-mass-transfer}
\textup{(Hamiltonian and Mass Transfer; O\_Mass.Matter Remainder.)}
\textup{[Conditional on Thm.~\ref{thm:vrp-energy-transfer}.]}
By Prop.~\ref{prop:vrp-transfer-reduction}, the branch Hamiltonian
and branch mass scale transfer functorially along the same lawful
interfaces.  The synthesis obligations O\_Ham.Transfer and
O\_Mass.Transfer, previously reduced to O\_Energy.Transfer, are
thereby conditionally discharged, and the remainder of
O\_Mass.Matter recorded in Rmk.~\ref{rmk:vrp-mass-matter} closes
at the same conditional force.
\end{corollary}

\begin{proof}
Immediate: Thm.~\ref{thm:vrp-energy-transfer} supplies the
hypothesis of Prop.~\ref{prop:vrp-transfer-reduction}. \qed
\end{proof}

\subsection{Time as Licensed Continuation Order}
\label{sec:vrp-time}

Time sits upstream of Energy in the recovery ladder; it is
promoted here, after the coercive rungs, without circularity:
nothing in the Energy, Hamiltonian, Mass, or Transfer material
above uses a time variable.  The promotion order is an editorial
fact, not a dependency.

\begin{definition}[\status{D}]\label{def:branch-time-order}
\textup{(Branch Continuation Order.)}
Let $B$ be a lawful NFC regime with observational quotient
$\Sigma_B$.  For quotient classes $x, y \in \Sigma_B$, write
$x \preceq_t y$ when $y$ is reachable from $x$ by a finite chain
of admissible continuation steps, each of which is a certified
quotient-visible state change --- a step whose occurrence and
persistence across a refinement stage is certified by the declared
admissible probes.  This is the branch instance of the Book~I
survivor-chain pattern (Def.~\ref{def:survivor-component}).  The
relation $\preceq_t$ is the \emph{branch continuation order}.
Branch time \emph{is} this order structure: an ordering of
certified change, not a parameter, coordinate, or clock.
\end{definition}

\begin{theorem}[\status{C}]\label{thm:vrp-time}
\textup{(VRP-Time: Time as Licensed Continuation Order.)}
\textup{[Conditional on: a lawful branch $B$ whose admissible
continuations form a stable system of certified quotient-visible
state changes (the stabilization discipline of Book~I,
Thm.~\ref{thm:stabilization}, supplies the spine instance); and
transfer-compatibility of lawful interfaces with continuation
(Def.~\ref{def:transfer-compatibility}).]}
The branch continuation order $\preceq_t$
(Def.~\ref{def:branch-time-order}) is a licensed physics variable
in the sense of \texttt{def:variable-licensing} (Book~VII),
realized through the quotient-visible-invariant disjunct.  The five
recovery clauses:
\begin{enumerate}[label=\textup{(\arabic*)}]
  \item \emph{Probe signature:} the finite admissible probes that
    certify a state change has occurred and persisted across a
    refinement step --- exactly the probes underlying survivor-chain
    membership.
  \item \emph{Quotient invariant:} $\preceq_t$ is a relation on
    quotient classes, defined through certified changes only, hence
    invariant under admissible re-description.
  \item \emph{Transport law:} lawful interfaces are monotone for
    $\preceq_t$: transfer-compatibility carries admissible
    continuation chains to admissible continuation chains, so
    $x \preceq_t y$ implies $I(x) \preceq_t I(y)$.
  \item \emph{Continuum encoding:} deferred --- the passage from
    the partial continuation order to a linear or smooth time
    parameter is exactly the open encoding obligation below; no
    $t \in \mathbb{R}$, $e^{-itH}$, or smooth-parameter notation
    is licensed by this theorem.
  \item \emph{Anti-smuggling audit:} no absolute time, manifold
    time coordinate, smooth parameter, or external clock is used;
    all content descends from refinement/continuation structure
    already certified in Books~I--III.
\end{enumerate}
\end{theorem}

\begin{proof}
Reflexivity and transitivity of $\preceq_t$ are immediate from
chain concatenation, so $\preceq_t$ is a preorder on the quotient.
Quotient-visibility (clause~2): each generating step is a certified
quotient-visible change, and the reachability relation built from
certified steps is itself determined by quotient data; an
admissible re-description fixes the quotient classes and the
admissibility of steps, hence fixes the relation.  Transport law
(clause~3): by Def.~\ref{def:transfer-compatibility} a lawful
interface carries admissible continuation structure to admissible
continuation structure, so the image of a generating chain is a
generating chain; monotonicity follows by induction on chain
length.  Spine instance: the Book~I survivor-chain order
(Def.~\ref{def:survivor-component}, Thm.~\ref{thm:stabilization})
is exactly the relation generated by certified survival steps, so
the spine carries the order natively; every branch's depth-indexed
machinery (depth-sum tails, $O(n^{-1})$ certificates, stabilization
depths) has been using this order implicitly throughout canon.
The declaration clause of the licensing predicate is met by the
itemization in the statement, and clause~(4)'s deferral is itself
part of the declaration: this theorem licenses the \emph{order},
not its numerical encoding. \qed
\end{proof}

\begin{remark}[\status{R}]\label{rmk:vrp-time-depth-index}
\textup{(Depth Index as De Facto Within-Chain Comparison; Clock
Comparison Narrowed.)}
Canon already operates a de facto comparison index: refinement
depth $n$ (depth-sum tails, $O(n^{-1})$ certificates, stabilization
depths).  This supplies \emph{within-chain} comparison of certified
changes.  The genuinely open content of the synthesis obligation
O\_Time.Clock is therefore \emph{cross-chain canonicity}: whether
equal certified step counts between \emph{different} continuation
chains constitute lawful synchronization or a presentation
artifact.  This remark narrows the obligation; it confers no
force.
\end{remark}

\begin{remark}[\status{C}]\label{ob:vrp-time-encode}
\textbf{O\_Time.Encode (Order-to-Number Encoding) ---
Conditionally Discharged; the residual re-registered, not
retained.}
\textup{[Conditional on Thm.~\ref{thm:vrp-entropy-clock} (clause
(a)), Thm.~\ref{thm:vrp-time-linear} and
Thm.~\ref{thm:vrp-time-smooth} (clause (b)),
\S\S\,\ref{sec:vrp-entropy-clock}, \ref{sec:vrp-time-linear}.]}
Clause~(a): the lawful clock is the entropy/defect ledger
(L+28).  Clause~(b): the discrete encoding is settled by
\emph{characterization} --- a canonical clock coordinate exists on
a sector iff the sector is $\Delta$-coherent, with affine gauge
uniqueness (unit and origin only) and a canonical linear
\emph{preorder} extending $\preceq_t$ without any Szpilrajn-style
choice; the continuum encoding holds as a bridge clause at the
same epistemic level as the LSC clause of the energy form.
What this discharge does \emph{not} license is a time parameter
for defect-free dynamics: that question is not an encoding
residual but a distinct problem, re-registered as
Rmk.~\ref{ob:vrp-time-reversible} (O\_Time.Reversible) rather
than silently retained here.
\end{remark}

\subsection{Evolution Generation: Monotone Realization of Branch
Time by the Licensed Generator}
\label{sec:vrp-ham-time}

This subsection discharges O\_Ham.Time
(Rmk.~\ref{ob:vrp-ham-time}).  It is the first theorem of the
corpus binding two licensed variables \emph{dynamically}: the
generator of the licensed energy form
(Thm.~\ref{thm:vrp-hamiltonian}) and the licensed continuation
order (Thm.~\ref{thm:vrp-time}).  The strategy is order-level
throughout: ``evolution'' means motion along $\preceq_t$, and no
smooth-time notation is used or licensed.

\begin{definition}[\status{D}]\label{def:hb-compatible-step}
\textup{($H_B$-Compatible Steps and the Generated Suborder.)}
Let $B$ carry a licensed energy form $\mathfrak{h}_B$
(Thm.~\ref{thm:vrp-energy}) with generator $H_B$
(Thm.~\ref{thm:vrp-hamiltonian}), a licensed continuation order
$\preceq_t$ (Thm.~\ref{thm:vrp-time}), and a declared dissipation
ledger on continuation steps.  An admissible continuation step
$x \to y$ is \emph{$H_B$-compatible} if
\begin{enumerate}[label=\textup{(\roman*)}]
  \item it is a generating step of $\preceq_t$ (admissible and a
    certified quotient-visible state change); and
  \item it is \emph{form-non-increasing}:
    $\mathfrak{h}_B(y) \le \mathfrak{h}_B(x)$, with the decrement
    $\mathfrak{h}_B(x) - \mathfrak{h}_B(y)$ recorded on the
    declared dissipation ledger.
\end{enumerate}
The \emph{$H_B$-generated suborder} $\preceq_H$ is the preorder
generated by $H_B$-compatible steps.  By construction
$\preceq_H \,\subseteq\, \preceq_t$.
\end{definition}

\begin{lemma}[\status{C}]\label{lem:vrp-ham-lyapunov}
\textup{(Energy as Lyapunov Ledger along $H_B$-Compatible
Continuation.)}
\textup{[Conditional on: Thm.~\ref{thm:vrp-energy},
Thm.~\ref{thm:vrp-hamiltonian}, and SCT.1
(\texttt{thm:SCT1}, Book~II) on the branch.]}
Along any $\preceq_H$-chain the coercive cost $\mathfrak{h}_B$ is
non-increasing, every decrement is ledger-visible, and the total
cost dissipated along the chain equals the sum of the per-step
ledger entries.  Moreover, the canonical contraction family of the
nonnegative generator is form-non-increasing: for every $s \ge 0$
(an internal spectral-calculus parameter, \emph{not} a licensed
time variable),
\[
  \overline{\mathfrak{h}_B}\bigl(e^{-sH_B}\psi\bigr)
  \;\le\;
  \overline{\mathfrak{h}_B}(\psi)
  \qquad
  \text{for all } \psi \in
  \mathrm{dom}\bigl(\overline{\mathfrak{h}_B}\bigr),
\]
so each application of the contraction family whose quotient image
is a certified state change is an $H_B$-compatible step.
\end{lemma}

\begin{proof}
Monotonicity along a $\preceq_H$-chain is
Def.~\ref{def:hb-compatible-step}(ii) applied stepwise; the
telescoping identity gives the total-dissipation accounting.
Ledger-visibility of every decrement is SCT.1: all transport
degradation is ledger-visible, and a decrement off the declared
ledger would be hidden loss, which SCT.1 excludes; the reverse
direction (no spontaneous coercive gain along admissible steps) is
transfer-coercivity (Def.~\ref{def:transfer-coercivity}).  For the
contraction family: $H_B \ge 0$ is self-adjoint
(Thm.~\ref{thm:vrp-hamiltonian}), so by the spectral theorem
$\overline{\mathfrak{h}_B}(e^{-sH_B}\psi) = \int_{[0,\infty)}
\lambda\, e^{-2s\lambda}\, d\|E_\lambda \psi\|^2 \le
\int_{[0,\infty)} \lambda\, d\|E_\lambda \psi\|^2 =
\overline{\mathfrak{h}_B}(\psi)$, since
$\lambda e^{-2s\lambda} \le \lambda$ on
$\operatorname{Spec}(H_B) \subseteq [0,\infty)$.  The parameter
$s$ enters only through functional calculus of the licensed
operator; no time-evolution claim is made or needed. \qed
\end{proof}

\begin{theorem}[\status{C}]\label{thm:vrp-ham-time-realization}
\textup{(Monotone Realization: When $H_B$-Dynamics Realizes Branch
Time.)}
\textup{[Conditional on: Thm.~\ref{thm:vrp-time},
Thm.~\ref{thm:vrp-hamiltonian},
Lem.~\ref{lem:vrp-ham-lyapunov}, SCT.1 on the branch, and the
declared dissipation-ledger discipline.]}
\begin{enumerate}[label=\textup{(\arabic*)}]
  \item \emph{(Order-compatibility.)}  The $H_B$-generated family
    is monotone for branch time: every $H_B$-compatible step is a
    generating step of $\preceq_t$, hence
    $\preceq_H \subseteq \preceq_t$ and motion under the
    $H_B$-generated family never contradicts the licensed
    continuation order.
  \item \emph{(Realization criterion.)}  $H_B$ \emph{realizes}
    branch time on a declared sector $S \subseteq \Sigma_B$ if and
    only if every generating step of $\preceq_t$ within $S$ is
    $H_B$-compatible up to ledger-visible defect --- equivalently,
    $\preceq_H|_S \,=\, \preceq_t|_S$.  In that case
    ``the evolution generated by $H_B$'' has licensed meaning on
    $S$ as motion along the branch continuation order, and the
    demand of O\_Ham.Time is met on $S$.
  \item \emph{(Failure diagnosis.)}  If
    $\preceq_H|_S \subsetneq \preceq_t|_S$, the discrepancy is
    ledger-visible: there exists a certified continuation step
    within $S$ whose coercive-cost change is not accounted by the
    $H_B$ dissipation ledger, witnessing either an energy channel
    outside the declared form (an undeclared hidden supplement,
    prohibited) or a continuation step imported from outside the
    admissible family (an imported flow).  The failure of
    realization is therefore always nameable, never silent.
\end{enumerate}
\emph{Scope clause:} this theorem operates at the order level
only.  $e^{-itH_B}$ with $t \in \mathbb{R}$ as licensed time
remains unlicensed pending O\_Time.Encode clause~(b)
(Rmk.~\ref{ob:vrp-time-encode}).
\end{theorem}

\begin{proof}
\textbf{(1)}  By Def.~\ref{def:hb-compatible-step}(i), every
$H_B$-compatible step is a generating step of $\preceq_t$;
monotonicity of the generated preorder follows by induction on
chain length, exactly as in the transport-law clause of
Thm.~\ref{thm:vrp-time}.  By Lem.~\ref{lem:vrp-ham-lyapunov} the
canonical contraction family of $H_B$ supplies form-non-increasing
steps, so the $H_B$-generated family moves only along
$\preceq_H \subseteq \preceq_t$.

\textbf{(2)}  ($\Leftarrow$)  Suppose every generating step of
$\preceq_t$ within $S$ is $H_B$-compatible up to ledger-visible
defect.  Then every generating chain of $\preceq_t|_S$ is a chain
of $H_B$-compatible steps with all defects ledgered, hence a
$\preceq_H$-chain; so $\preceq_t|_S \subseteq \preceq_H|_S$, and
with~(1) the two preorders coincide on $S$.  ($\Rightarrow$)
Conversely, if $\preceq_H|_S = \preceq_t|_S$ then in particular
each generating step $x \to y$ of $\preceq_t$ within $S$ satisfies
$x \preceq_H y$, i.e.\ it is reachable by $H_B$-compatible steps
with ledgered accounting --- which is the criterion.  Under the
criterion, ``evolution'' is motion along $\preceq_t$ realized by
$H_B$-compatible steps; since $\preceq_t$ is the licensed time
variable (Thm.~\ref{thm:vrp-time}) and $H_B$ the licensed
generator (Thm.~\ref{thm:vrp-hamiltonian}), the phrase
``evolution generated by $H_B$'' carries exactly the licensed
content and no more.

\textbf{(3)}  Suppose $x \preceq_t y$ within $S$ but
$x \not\preceq_H y$.  Every generating chain from $x$ to $y$ then
contains a step $x' \to y'$ that fails
Def.~\ref{def:hb-compatible-step}(ii) modulo the ledger: either
$\mathfrak{h}_B(y') > \mathfrak{h}_B(x')$ with the increment
unledgered, or the decrement disagrees with the ledger entry.  By
SCT.1 all lawful transport degradation is ledger-visible, so an
unledgered coercive change cannot arise from the declared
transport-coercion structure; it therefore witnesses either
coercive content entering through a channel outside the declared
form (a hidden supplement, prohibited by the standing
no-hidden-supplementation discipline) or a step that is not
generated by the admissible continuation family at all (an
imported flow).  Both diagnoses are recorded objects: the
discrepancy set is ledger-visible by construction. \qed
\end{proof}

\begin{remark}[\status{R}]\label{rmk:vrp-ham-time-instances}
\textup{(Realization Instances: NS Dissipative Discipline, YM
Phase-A Persistence.)}
Two certified canonical regimes satisfy the realization criterion
of Thm.~\ref{thm:vrp-ham-time-realization}(2) on their declared
sectors.  \emph{NS:} on the safe-tail regime the contraction
$\mathfrak{D}_{n+1} \le \kappa\,\mathfrak{D}_n$
(\texttt{cor:NS61-discharged}, NS Branch) makes every admissible
continuation step strictly form-non-increasing with the decrement
carried on the obstruction ledger, so every $\preceq_t$-generating
step is $H_B$-compatible and NS realizes branch time on the
declared sector.  \emph{YM:} on the Phase-A screened sector the
depth-indexed persistence machinery
(\texttt{thm:O-GLOB}, YM Branch) certifies each refinement stage
as a quotient-visible change while the contraction family of the
Friedrichs generator of \texttt{thm:B1-closable} is
form-non-increasing (Lem.~\ref{lem:vrp-ham-lyapunov}); the
Phase-A sector is a realization sector.  These are realization
instances of an already-proved criterion; this remark confers no
new force.
\end{remark}

\subsection{The Entropy Clock: Canonical Cross-Chain Comparison}
\label{sec:vrp-entropy-clock}

The promotion of the foundation rung pays an immediate dividend on
the temporal rung.  Rmk.~\ref{rmk:vrp-time-depth-index} narrowed
O\_Time.Clock to cross-chain canonicity of step counts; the Book~I
ledger machinery shows that raw step count was the wrong clock
candidate all along, and supplies the right one.

\begin{theorem}[\status{C}]\label{thm:vrp-entropy-clock}
\textup{(Entropy Clock: Canonical Cross-Chain Comparison via the
Defect Ledger.)}
\textup{[Conditional on: Thm.~\ref{thm:vrp-entropy},
Thm.~\ref{thm:vrp-time}; the Book~I duality and preorder
machinery (Thm.~\ref{thm:ledger-entropy-duality},
Prop.~\ref{prop:ledger-additivity}, Prop.~\ref{prop:defect-nonneg})
for the spine instance; and the branch carrying a declared defect
ledger $\Delta_B$ compatible with its admissible continuation
steps.]}
\begin{enumerate}[label=\textup{(\arabic*)}]
  \item \emph{(Canonicity.)}  Cumulative defect weight $\Delta_B$
    is a quotient-visible, reparametrization-invariant chain
    functional: its value along a continuation chain depends only
    on the certified class-splitting and coarsening events the
    chain traverses, not on the presentation or step-partition of
    the chain.  By Ledger--Entropy Duality, the cumulative defect
    along a chain equals the entropy decrement, so the entropy
    ledger and the defect ledger define the \emph{same} clock.
  \item \emph{(Order-compatibility.)}  The $\Delta_B$-clock is
    monotone for branch time: along any admissible continuation
    $\Delta_B$ is non-decreasing, so $x \preceq_t y$ implies that
    every generating chain from $x$ to $y$ carries
    $\Delta_B$-weight $\ge 0$, and the clock never runs backward
    against the licensed order.
  \item \emph{(Cross-chain synchronization; O\_Time.Encode
    clause~(a) conditionally discharged.)}  Two continuation
    chains are lawfully synchronized exactly when their cumulative
    defect weights agree at corresponding certified events.  Equal
    raw step count is neither necessary nor sufficient: chain
    depth is a lawful clock precisely on families of chains with
    uniform per-step defect, and otherwise is a presentation
    index.  The lawful clock demanded by O\_Time.Encode
    clause~(a) is the entropy/defect ledger.
\end{enumerate}
\end{theorem}

\begin{proof}
\textbf{(1)}  Each entry of $\Delta_B$ records a certified
class-splitting or coarsening event --- a quotient-visible datum
by construction (the defect ledger is the transaction log of the
quotient, Def.~\ref{def:defect-ledger}, Book~I).
Re-partitioning a chain into different steps re-groups the event
multiset but does not alter it; by ledger additivity
(Prop.~\ref{prop:ledger-additivity}) the cumulative weight is the
sum over the event multiset, hence partition-independent and
re-description-invariant.  The identity with the entropy decrement
is Ledger--Entropy Duality
(Thm.~\ref{thm:ledger-entropy-duality}) on spine survivor chains,
carried to branch continuation chains by clause~\ref{es:ledger} of
the licensed branch ledger (Thm.~\ref{thm:vrp-entropy}).

\textbf{(2)}  Each admissible continuation step appends a
nonnegative one-step defect (Prop.~\ref{prop:defect-nonneg});
induction on chain length gives monotonicity along every
generating chain of $\preceq_t$, which is order-compatibility in
the same sense as Thm.~\ref{thm:vrp-ham-time-realization}(1).

\textbf{(3)}  ($\Leftarrow$)  If two chains agree in cumulative
defect at corresponding certified events, then by~(1) they agree
in every quotient-visible clock reading, and by duality in entropy
decrement; no admissible probe can distinguish their temporal
content, which is lawful synchronization.  ($\Rightarrow$)
Conversely, a disagreement in cumulative defect at corresponding
events is itself a certified, quotient-visible discrepancy, so the
chains are not synchronized.  For step count: on a family of
chains with uniform per-step defect $\delta_0$, depth $n$ and
weight $n\delta_0$ are affinely equivalent, so depth is lawful
there; in general two chains of equal depth may traverse different
event multisets, and~(1) shows the weight, not the count, is the
invariant.  This is exactly the canonicity demanded by
O\_Time.Encode clause~(a), conditionally discharged with the
clock identified. \qed
\end{proof}

\begin{remark}[\status{R}]\label{rmk:vrp-entropy-clock-arrow}
\textup{(The Clock and the Arrow Are One Object.)}
The Book~I irreversible temporal preorder
(Sch.~\ref{sch:temporal-preorder}: $h \preceq h'$ iff
$\Delta(h) \le \Delta(h')$, forward-only by defect nonnegativity)
is, after this subsection, recognizable as the spine instance of
the entropy clock read as an order: the arrow of time and the
clock of time are the same ledger viewed as order and as scale.
No new force; this remark records the identification for the
encoding work that remains.
\end{remark}

\subsection{The Canonical Nontrivial Sector: O\_Mass.Sector
Discharged by Three-Variable Binding}
\label{sec:vrp-mass-sector}

The sector hypothesis of Thm.~\ref{thm:vrp-mass} is the last piece
of that theorem still entering as branch-supplied data.  This
subsection recovers it.  The construction binds \emph{three}
licensed variables --- the entropy ledger, the continuation order,
and the generator --- and its shape is forced by an audit failure
recorded inside the theorem itself.

\begin{definition}[\status{D}]\label{def:vrp-nontrivial-sector}
\textup{(Persistent, Label, and Nontrivial Sectors.)}
Let $B$ carry the universal entropy ledger $S_B$
(Thm.~\ref{thm:vrp-entropy-univ}), the licensed continuation order
$\preceq_t$ (Thm.~\ref{thm:vrp-time}), and the licensed generator
$H_B$ in its declared Hilbert realization
(Thm.~\ref{thm:vrp-hamiltonian}).  Define:
\begin{enumerate}[label=\textup{(\roman*)}]
  \item the \emph{persistent-distinction classes}
    \[
      \Sigma_B^{\mathrm{pers}} \;:=\;
      \{\, x \in \Sigma_B \;:\; S_B(y) > 0
      \text{ for every } y \succeq_t x \,\},
    \]
    the classes from which every admissible future sustains at
    least one admissible distinction (note
    $\Sigma_B^{\mathrm{pers}}$ is up-closed for $\preceq_t$ by
    construction);
  \item the \emph{persistent sector}
    $\mathcal{H}_B^{\mathrm{pers}}$: the closed span, in the
    declared realization, of the states realizing classes of
    $\Sigma_B^{\mathrm{pers}}$;
  \item the \emph{label sector}
    $\mathcal{H}_B^{\mathrm{label}} :=
    \mathcal{H}_B^{\mathrm{pers}} \cap \ker H_B$: persistent
    distinctions of zero coercive cost;
  \item the \emph{canonical nontrivial sector}
    \[
      \mathcal{H}_B^{\mathrm{nontriv}} \;:=\;
      \mathcal{H}_B^{\mathrm{pers}} \ominus
      \mathcal{H}_B^{\mathrm{label}},
    \]
    persistent distinctions of positive coercive cost.
\end{enumerate}
Here $\ker H_B = \{\psi : \overline{\mathfrak{h}_B}(\psi) = 0\}$
is the kernel of the licensed closed form (for a nonnegative
self-adjoint generator, $\overline{\mathfrak{h}_B}(\psi) =
\|H_B^{1/2}\psi\|^2$, so the two kernels coincide).
\end{definition}

\begin{theorem}[\status{C}]\label{thm:vrp-mass-sector}
\textup{(O\_Mass.Sector Discharge: the Canonical Nontrivial
Sector.)}
\textup{[Conditional on: Thm.~\ref{thm:vrp-entropy-univ},
Thm.~\ref{thm:vrp-time}, Thm.~\ref{thm:vrp-hamiltonian}; and
sector--realization compatibility: the declared realization
carries the $\preceq_t$-forward-invariance of
$\Sigma_B^{\mathrm{pers}}$ to invariance of
$\mathcal{H}_B^{\mathrm{pers}}$ under the contraction family of
$H_B$ --- automatic on realization sectors in the sense of
Thm.~\ref{thm:vrp-ham-time-realization}(2).]}
\begin{enumerate}[label=\textup{(\arabic*)}]
  \item \emph{(Quotient-visibility; no branch-supplied data.)}
    $\mathcal{H}_B^{\mathrm{nontriv}}$ is determined by the three
    licensed variables and the declared realization alone:
    admissible re-description fixes $S_B$, $\preceq_t$, and
    $\overline{\mathfrak{h}_B}$ (hence $\ker H_B$), and therefore
    fixes the sector.  No sector datum is imported.
  \item \emph{(Reduction; the spectral threshold is
    well-posed.)}  $\mathcal{H}_B^{\mathrm{pers}}$ reduces $H_B$,
    hence so does $\mathcal{H}_B^{\mathrm{nontriv}}$, and
    $\inf \operatorname{Spec}\bigl(H_B|_{\mathcal{H}_B^{\mathrm{nontriv}}}\bigr)$
    is well-defined.  Substituting this sector into
    Thm.~\ref{thm:vrp-mass} turns its sector hypothesis into
    recovered structure; only the \emph{positivity} of the
    threshold remains branch data, by design.
  \item \emph{(Coercivity audit, passed by forced
    factorization.)}  The naive candidate
    $\mathcal{H}_B^{\mathrm{pers}}$ alone fails the audit: it
    contains persistent distinctions of zero coercive cost, on
    which the spectral threshold is $0$, conflating
    distinguishability with excitation --- the same gap that
    blocked the entropy-to-energy reduction
    (Rmk.~\ref{rmk:vrp-univ-method-scope}).  The factorization
    through $\ker H_B$ is therefore forced; it is licensed because
    the kernel of a licensed closed form is canonical.
  \item \emph{(Byproduct: the label sector.)}
    $\mathcal{H}_B^{\mathrm{label}}$ is a derived canonical
    object: persistent, distinguishable, costless labels, stable
    under the $H_B$-generated family.  Its phenomenology matches
    superselection-type structure; the identification with branch
    superselection sectors is recorded as an observation, not
    claimed as a theorem.
\end{enumerate}
\end{theorem}

\begin{proof}
\textbf{(1)}  Each ingredient is quotient-determined:
$S_B$ by Thm.~\ref{thm:vrp-entropy-univ}(2), $\preceq_t$ by
Thm.~\ref{thm:vrp-time} clause~(2), and
$\overline{\mathfrak{h}_B}$ by Thm.~\ref{thm:vrp-energy} clause~(2)
with its kernel determined by the form.  The set
$\Sigma_B^{\mathrm{pers}}$ is defined by a condition in $S_B$ and
$\preceq_t$ only; its span, the kernel intersection, and the
orthogonal complement are determined by these data within the
declared realization.  An admissible re-description fixes all
three variables, hence the sector.

\textbf{(2)}  $\Sigma_B^{\mathrm{pers}}$ is up-closed for
$\preceq_t$: if every admissible future of $x$ has positive
burden, the same holds for every $y \succeq_t x$, whose futures
are among the futures of $x$.  By
Thm.~\ref{thm:vrp-ham-time-realization}(1) the $H_B$-generated
family moves only along $\preceq_t$, so at the quotient level it
preserves $\Sigma_B^{\mathrm{pers}}$; by the sector--realization
compatibility hypothesis this lifts to
$e^{-sH_B}\,\mathcal{H}_B^{\mathrm{pers}} \subseteq
\mathcal{H}_B^{\mathrm{pers}}$ for all $s \ge 0$.  Since each
$e^{-sH_B}$ is self-adjoint, invariance of the closed subspace
$\mathcal{H}_B^{\mathrm{pers}}$ implies invariance of its
orthogonal complement, so $\mathcal{H}_B^{\mathrm{pers}}$ reduces
every $e^{-sH_B}$ and hence reduces $H_B$.  The kernel
$\ker H_B$ is the fixed-point space of the contraction family and
reduces $H_B$ trivially; intersections and orthogonal differences
of reducing subspaces reduce, so
$\mathcal{H}_B^{\mathrm{nontriv}}$ reduces $H_B$ and the
restricted spectrum is well-defined.  The substitution claim is
immediate: Thm.~\ref{thm:vrp-mass} required a quotient-visible
sector; clause~(1) supplies it with no branch input, leaving
$m_B^2 > 0$ as the sole branch-data hypothesis --- which that
theorem already declared as such.

\textbf{(3)}  On $\mathcal{H}_B^{\mathrm{label}} =
\mathcal{H}_B^{\mathrm{pers}} \cap \ker H_B$ the form vanishes,
so $0 \in \operatorname{Spec}(H_B|_{\mathcal{H}_B^{\mathrm{pers}}})$
whenever the label sector is nonzero, and the threshold on the
unfactored candidate is $0$ regardless of any gap above the
kernel.  Distinguishability ($S_B > 0$, class-counting content)
does not imply coercive cost ($\mathfrak{h}_B > 0$, transport
content); the audit gap is exactly the one isolated in
Rmk.~\ref{rmk:vrp-univ-method-scope}, and the factorization
removes precisely the zero-cost part and nothing else.

\textbf{(4)}  Stability of the label sector under the
$H_B$-generated family holds because $\ker H_B$ is the fixed-point
space and $\mathcal{H}_B^{\mathrm{pers}}$ is invariant by~(2);
persistence and distinguishability hold by membership in
$\Sigma_B^{\mathrm{pers}}$.  No further claim is made. \qed
\end{proof}

\begin{remark}[\status{R}]\label{rmk:vrp-mass-sector-instances}
\textup{(Instance Verification: YM Phase-A.)}
The construction recovers the certified instance: on the YM
Phase-A screened sector, persistence is supplied by the
depth-indexed persistence machinery (\texttt{thm:O-GLOB}, YM
Branch), the screening removes the zero-cost part, and the
positive threshold is the certified gap
$\Delta_{\mathrm{YM}} = m^2 > 0$ --- so the Phase-A screened
sector instantiates
$\mathcal{H}_B^{\mathrm{nontriv}}$ with the realization-sector
hypothesis certified at L+27
(Rmk.~\ref{rmk:vrp-ham-time-instances}).  NS realizes the
construction trivially in the complementary sense: its declared
regime concerns dissipation rather than a mass gap, and no gap
claim is made.  This remark records instances; no new force.
\end{remark}

\subsection{The Canonical Domain: O\_Ham.Domain Discharged by
Contraction Regularization}
\label{sec:vrp-ham-domain}

The last operator-theoretic gap in the Hamiltonian rung is the
domain: Thm.~\ref{thm:vrp-hamiltonian} supplies
$\mathrm{dom}(\overline{\mathfrak{h}_B})$ as the form-norm
completion of the declared core, but the obligation demands that
this domain be \emph{recovered} from finite probe and refinement
data rather than presentation-supplied.  The naive statement ---
``the domain is the closure of the finite-probe core'' --- would be
a tautology if the core property were hypothesized; the content of
this subsection is that the core property can be \emph{derived},
because the licensed contraction family regularizes finite-probe
data into a form core.

\begin{definition}[\status{D}]\label{def:vrp-finite-probe-core}
\textup{(Finite-Probe States; the Regularized Core.)}
Let $B$ satisfy the hypotheses of Thm.~\ref{thm:vrp-hamiltonian}
with declared Hilbert realization $\mathcal{H}_B$.  A
\emph{finite-probe state} is a state of $\mathcal{H}_B$ realizing
a quotient class certified at finite refinement depth; write
$\mathcal{D}_B^{\mathrm{fin}}$ for their linear span.  The
\emph{contraction-regularized finite-probe core} is
\[
  \mathcal{D}_B^{\mathrm{reg}} \;:=\;
  \mathrm{span}\bigl\{\, e^{-sH_B}\psi \;:\;
  \psi \in \mathcal{D}_B^{\mathrm{fin}},\; s \ge 0 \,\bigr\},
\]
where the contraction family is the licensed functional calculus
of the generator (Lem.~\ref{lem:vrp-ham-lyapunov}: $s$ is an
internal spectral parameter, not a licensed time variable).
\end{definition}

\begin{theorem}[\status{C}]\label{thm:vrp-ham-domain}
\textup{(O\_Ham.Domain Discharge: the Domain as Recovered
Structure.)}
\textup{[Conditional on: Thm.~\ref{thm:vrp-hamiltonian};
\emph{realization faithfulness} --- the finite-probe states are
norm-dense in $\mathcal{H}_B$, which is the Hilbert-level
expression of the no-hidden-channel discipline
(\texttt{def:hidden-channel}, Book~V): a state orthogonal to
every finite-depth-certified state is invisible to every
admissible probe, i.e.\ an undeclared channel of the realization;
and \emph{finite per-stage cost} --- every finite-probe state lies
in $\mathrm{dom}(\mathfrak{h}_B)$.]}
\begin{enumerate}[label=\textup{(\arabic*)}]
  \item \emph{(Canonical core.)}
    $\mathcal{D}_B^{\mathrm{reg}}$ is a form core:
    $\mathrm{dom}(\overline{\mathfrak{h}_B})$ is the closure of
    $\mathcal{D}_B^{\mathrm{reg}}$ in the form norm
    $\|\psi\|_{\mathfrak{h}}^2 = \|\psi\|^2 +
    \overline{\mathfrak{h}_B}(\psi)$.
  \item \emph{(Quotient-visible membership criterion.)}
    $\psi \in \mathrm{dom}(\overline{\mathfrak{h}_B})$ if and only
    if $\psi$ is the norm limit of a form-Cauchy sequence from
    $\mathcal{D}_B^{\mathrm{reg}}$ --- a criterion expressed
    entirely in licensed data: finite-probe certificates, the
    licensed contraction family, and the licensed form.  The
    domain is recovered structure, not presentation-supplied.
  \item \emph{(Presentation independence.)}  Any two admissible
    finite-probe presentations that are mutually certifiable
    (each state of one is norm-approximable by certified states
    of the other --- the admissible re-description discipline)
    yield regularized cores with the same form-norm closure.
  \item \emph{(Finality.)}  Closedness of
    $\overline{\mathfrak{h}_B}$ makes the characterization final:
    no admissible refinement enlarges the domain beyond this
    closure, matching the discharge-by-construction of O\_Ham.SA
    in Thm.~\ref{thm:vrp-hamiltonian}.
\end{enumerate}
Anti-smuggling audit: no Sobolev scale, manifold smooth
structure, or a-priori operator domain is imported; the domain
descends from probe certificates and licensed functional calculus
alone.
\end{theorem}

\begin{proof}
\textbf{(1)}  \emph{Density in $\mathcal{H}_B$:} by strong
continuity of the contraction family of a nonnegative
self-adjoint generator, $e^{-sH_B}\psi \to \psi$ as
$s \to 0^+$; with realization faithfulness,
$\mathcal{D}_B^{\mathrm{reg}} \supseteq
\mathcal{D}_B^{\mathrm{fin}}$ is norm-dense.
\emph{Containment:} for $s > 0$, $e^{-sH_B}\psi \in
\mathrm{dom}(H_B^k)$ for every $k$; for $s = 0$ the finite
per-stage cost hypothesis gives
$\psi \in \mathrm{dom}(\mathfrak{h}_B)$.  Hence
$\mathcal{D}_B^{\mathrm{reg}} \subseteq
\mathrm{dom}(\overline{\mathfrak{h}_B})$.
\emph{Invariance:} $e^{-tH_B}e^{-sH_B}\psi =
e^{-(s+t)H_B}\psi \in \mathcal{D}_B^{\mathrm{reg}}$.

It remains to show every
$\psi \in \mathrm{dom}(\overline{\mathfrak{h}_B})$ is
form-approximable from $\mathcal{D}_B^{\mathrm{reg}}$.  Two
spectral estimates suffice.  \emph{First}, for
$\psi \in \mathrm{dom}(\overline{\mathfrak{h}_B})$,
\[
  \|e^{-sH_B}\psi - \psi\|_{\mathfrak{h}}^2
  \;=\; \int_{[0,\infty)} (1+\lambda)\,
  \bigl|e^{-s\lambda} - 1\bigr|^2 \, d\|E_\lambda\psi\|^2
  \;\xrightarrow[s\to 0^+]{}\; 0
\]
by dominated convergence: the integrand is dominated by
$4(1+\lambda)$, integrable since
$\psi \in \mathrm{dom}(\overline{\mathfrak{h}_B})$, and tends to
$0$ pointwise.  So it suffices to form-approximate
$e^{-sH_B}\psi$ for fixed $s > 0$.  \emph{Second}, the smoothing
bound: for fixed $s > 0$ and any $\phi \in \mathcal{H}_B$,
\[
  \|e^{-sH_B}\phi\|_{\mathfrak{h}}^2
  \;=\; \int_{[0,\infty)} (1+\lambda)\, e^{-2s\lambda}\,
  d\|E_\lambda\phi\|^2
  \;\le\; C_s\, \|\phi\|^2,
  \qquad
  C_s := \sup_{\lambda \ge 0}\,(1+\lambda)e^{-2s\lambda} < \infty,
\]
i.e.\ the contraction family is bounded from the Hilbert norm to
the form norm.  Now take finite-probe $\psi_n \to \psi$ in norm
(faithfulness); then $e^{-sH_B}\psi_n \in
\mathcal{D}_B^{\mathrm{reg}}$ and
$\|e^{-sH_B}(\psi_n - \psi)\|_{\mathfrak{h}}^2 \le
C_s \|\psi_n - \psi\|^2 \to 0$.  Composing the two
approximations gives a form-norm approximation of $\psi$ from
$\mathcal{D}_B^{\mathrm{reg}}$, which is the core property; the
form-norm closure of a core is
$\mathrm{dom}(\overline{\mathfrak{h}_B})$ by definition of the
closure.

\textbf{(2)}  Restates~(1) through the definition of form-norm
closure.  Every ingredient of the criterion is licensed: the
certificates are finite refinement data, the contraction family is
licensed functional calculus of the licensed generator
(Lem.~\ref{lem:vrp-ham-lyapunov}), and the form is licensed
(Thm.~\ref{thm:vrp-energy}).

\textbf{(3)}  Mutual certifiability gives mutual norm-density of
the two finite-probe spans in each other's closures, hence both
are norm-dense in $\mathcal{H}_B$; the smoothing bound then shows
each regularized core form-approximates the other, so both
closures equal $\mathrm{dom}(\overline{\mathfrak{h}_B})$ by~(1).

\textbf{(4)}  The closure of a closable form is closed, and the
Friedrichs domain is uniquely determined by the form
(uniqueness clause of Thm.~\ref{thm:vrp-hamiltonian}); any
admissible refinement acts within the same licensed form and
cannot pass beyond its closure. \qed
\end{proof}

\begin{remark}[\status{R}]\label{rmk:vrp-ham-domain-faithfulness}
\textup{(Faithfulness Is Not an Import.)}
The realization-faithfulness hypothesis deserves its own record:
it is not an analytic convenience but the Hilbert-level face of
Book~V lawfulness.  A nonzero vector orthogonal to every
finite-probe state would carry observational content reachable by
no admissible probe at any finite depth --- precisely an
undeclared dependence of the realization on structure invisible
through the declared projection, i.e.\ a hidden channel
(\texttt{def:hidden-channel}, Book~V).  Lawful branches exclude
it by discipline, not by assumption of convenience.  No force
conferred.
\end{remark}

\subsection{Linearization of Branch Time: O\_Time.Encode
Clause~(b) by Characterization, and the Reversibility Obstruction}
\label{sec:vrp-time-linear}

The clock theorem (\S\,\ref{sec:vrp-entropy-clock}) identified the
canonical cross-chain comparison; what remained of
O\_Time.Encode was the order-to-number encoding itself.  This
subsection settles the discrete encoding by a full
characterization (not merely a sectoral instance), supplies the
continuum encoding as a bridge clause at the same epistemic level
as the LSC clause of the energy form, and isolates --- as a new,
sharply named obligation --- the one question the entropy clock
provably cannot answer.

\begin{definition}[\status{D}]\label{def:vrp-coherent-sector}
\textup{($\Delta$-Coherence; Clock Coordinate; Instants;
Clock-Simultaneity.)}
Let $B$ carry the licensed continuation order $\preceq_t$
(Thm.~\ref{thm:vrp-time}) and the canonical clock
(Thm.~\ref{thm:vrp-entropy-clock}).  A $\preceq_t$-connected
subset $T \subseteq \Sigma_B$ with a declared origin $o$ is
\emph{$\Delta$-coherent} if any two generating chains between the
same endpoints within $T$ carry equal cumulative defect weight.
On a coherent sector define the \emph{clock coordinate}
\[
  \tau(x) \;:=\; \Delta_B(\text{any generating chain }
  o \leadsto x),
\]
well-defined by coherence.  Call $x \sim y$
(\emph{instant-equivalent}) when $x$ and $y$ are connected by a
chain of zero-defect steps, and call $x, y$
\emph{clock-simultaneous} when $\tau(x) = \tau(y)$.  A
\emph{uniform-defect family} is a family of chains of constant
per-step defect $\delta_0 > 0$.
\end{definition}

\begin{theorem}[\status{C}]\label{thm:vrp-time-linear}
\textup{(Linearization by Characterization: the Discrete Encoding
of Branch Time.)}
\textup{[Conditional on: Thm.~\ref{thm:vrp-time},
Thm.~\ref{thm:vrp-entropy-clock}.]}
\begin{enumerate}[label=\textup{(\arabic*)}]
  \item \emph{(Characterization.)}  A canonical clock coordinate
    (a quotient-visible state function assigning every generating
    step its ledgered weight) exists on a
    $\preceq_t$-connected sector $T$ \emph{if and only if} $T$ is
    $\Delta$-coherent.  Necessity is immediate: path-dependent
    weight is inconsistent with any state function reproducing
    per-step ledger entries; sufficiency is
    Def.~\ref{def:vrp-coherent-sector}.
  \item \emph{(Affine gauge uniqueness.)}  On a coherent sector,
    $\tau$ is monotone for $\preceq_t$, quotient-visible, and
    unique up to the clock gauge: any monotone encoding $\tau'$
    assigning each generating step its ledgered weight in a
    declared unit satisfies $\tau' = a\,\tau + b$ ($a$ the unit
    ratio, $b$ the origin shift), by telescoping along any chain
    and coherence.  The licensed time coordinate carries exactly
    the gauge freedom of a physical clock --- unit and origin ---
    and no more.
  \item \emph{(Strict monotonicity modulo instants; the fibers.)}
    For comparable $x \prec_t y$: $\tau(x) = \tau(y)$ iff $x \sim
    y$ (every connecting step is zero-defect), and otherwise
    $\tau(x) < \tau(y)$.  The fibers of $\tau$ consist of
    instant-classes together with $\preceq_t$-incomparable
    clock-simultaneous elements.
  \item \emph{(Canonical linear preorder; Szpilrajn avoided.)}
    The relation $x \le_\tau y :\iff \tau(x) \le \tau(y)$ is a
    \emph{linear preorder} extending $\preceq_t$ on $T$,
    determined by licensed data plus gauge --- no extension choice
    is made.  A Szpilrajn-style linear \emph{order} refining
    incomparable clock-simultaneous pairs would require exactly
    such a choice; canonicity is preserved by declining it:
    licensed time totally orders clock readings and deliberately
    does not refine simultaneity.  Anti-smuggling:
    clock-simultaneity is a bookkeeping notion (equal ledgered
    weight), not a causal, metric, or signal-theoretic claim;
    cross-sector comparison remains governed by
    Thm.~\ref{thm:vrp-entropy-clock}(3).
  \item \emph{(Uniform-defect specialization.)}  On a
    uniform-defect family, $\tau = \delta_0 \cdot
    \mathrm{depth}$: chain depth is an affine clock, recovering
    the criterion of Thm.~\ref{thm:vrp-entropy-clock}(3) as the
    special case.
\end{enumerate}
\end{theorem}

\begin{proof}
\textbf{(1)}  If a state function $\tau'$ reproduces per-step
ledger entries, then for any chain $o \leadsto x$ the telescoped
sum gives $\tau'(x) - \tau'(o) = \Delta_B(\text{chain})$, so all
chains to $x$ carry equal weight: coherence is necessary.
Conversely, on a coherent sector
Def.~\ref{def:vrp-coherent-sector} defines $\tau$ consistently.

\textbf{(2)}  Monotonicity is defect-nonnegativity along chains
(Thm.~\ref{thm:vrp-entropy-clock}(2)); quotient-visibility is
Thm.~\ref{thm:vrp-entropy-clock}(1).  For $\tau'$ as hypothesized,
telescoping along any chain $o \leadsto x$ gives
$\tau'(x) = \tau'(o) + a\,\Delta_B(\text{chain}) = b + a\,\tau(x)$,
chain-independent by coherence.

\textbf{(3)}  For $x \prec_t y$, every connecting chain has
nonnegative per-step weights summing to $\tau(y) - \tau(x)$; the
sum vanishes iff every step is zero-defect, i.e.\ $x \sim y$.
The fiber description follows.

\textbf{(4)}  $\le_\tau$ is reflexive, transitive, and total
(values in $\mathbb{R}$), and extends $\preceq_t$ by~(2)--(3).
It is determined by $\tau$, hence by licensed data and the gauge
of~(2).  That no further canonical refinement exists is the
content of~(3): an incomparable clock-simultaneous pair is
separated by no licensed datum, so any ordering of it is a choice
--- which the linear preorder declines by construction.

\textbf{(5)}  Immediate: constant per-step defect $\delta_0$
makes every chain weight $\delta_0 n$ at depth $n$. \qed
\end{proof}

\begin{theorem}[\status{C}]\label{thm:vrp-time-smooth}
\textup{(Continuum Encoding: the Bridge Clause.)}
\textup{[Conditional on: Thm.~\ref{thm:vrp-time-linear} on each
sector; and a declared Book~VI bridge presentation of the clock:
a refining family $\{T_k\}$ of $\Delta$-coherent sectors with
admissible refinement maps under which the clock coordinates form
a compatible monotone family, the value sets become
$\varepsilon_k$-dense in a declared interval with
$\varepsilon_k \to 0$, and the refinement defect of clock
readings vanishes in the limit (the clock analogue of the LSC
clause~\ref{en:lsc}).]}
There is a unique monotone continuous limit coordinate
$t$ on the bridge sector restricting compatibly to every
$\tau_k$ up to gauge; continuum time notation $t \in \mathbb{R}$
is thereby licensed on the bridge sector \emph{for order and
irreversible content}, with regularity beyond continuity
inherited from, and only from, the declared bridge regularity.
\end{theorem}

\begin{proof}
Compatibility and monotonicity give a well-defined monotone
function on the union of value sets; $\varepsilon_k$-density with
$\varepsilon_k \to 0$ makes this union dense in the declared
interval, and a monotone function on a dense set with vanishing
refinement defect extends uniquely to a continuous monotone
function (the vanishing defect supplies the modulus of
continuity, exactly as the LSC clause controls the energy form
across the bridge).  Uniqueness up to gauge follows from
Thm.~\ref{thm:vrp-time-linear}(2) applied at every stage.  The
scope restriction is part of the statement: nothing here
licenses a parameter for defect-free dynamics --- see
Rmk.~\ref{ob:vrp-time-reversible}. \qed
\end{proof}

\begin{remark}[\status{C}]\label{ob:vrp-time-reversible}
\textbf{O\_Time.Reversible --- Conditionally Discharged by
Dichotomy; the residual re-registered.}
\textup{[Conditional on
Thm.~\ref{thm:vrp-reversible-dichotomy},
Cor.~\ref{cor:vrp-no-global-unitary-time},
\S\,\ref{sec:vrp-reversible}.]}
The question as registered --- whether a second clock for
defect-free dynamics exists --- is answered by the dichotomy:
on the gauge horn no clock exists because no licensed change
exists ($e^{-itH_B}$ is admissible re-description); on the phase
horn the clock exists and is constructed --- cyclic/toral phase
from licensed spectral differences, canonical up to
origin~$+$~orientation gauge.  On both horns, global linear
unitary time is unlicensed, so thermodynamic and unitary time
remain provably distinct.  The genuinely open residual ---
canonical cross-orbit phase alignment, the defect-free analogue
of clock comparison --- is re-registered as
Rmk.~\ref{ob:vrp-phase-sync} (O\_Phase.Sync), not retained here.
\end{remark}

\subsection{The Reversible-Clock Dichotomy: O\_Time.Reversible
Resolved}
\label{sec:vrp-reversible}

The L+32 obstruction asked whether any licensed datum can
distinguish points of a unitary orbit --- the prerequisite for any
clock on defect-free dynamics.  The feasibility analysis resolves
into a dichotomy decidable from the declared probe family, with
both horns in-toolkit; this subsection proves it, constructs the
second clock on the horn where it exists, and shows that on
\emph{both} horns a global linear unitary time remains unlicensed
--- so the question splits correctly rather than closing falsely.

\begin{theorem}[\status{C}]\label{thm:vrp-reversible-dichotomy}
\textup{(The Reversible-Clock Dichotomy.)}
\textup{[Conditional on: Thm.~\ref{thm:vrp-hamiltonian},
Thm.~\ref{thm:vrp-time-linear}; the declared admissible probe
family of $B$; and, for the phase-coordinate clause, orbit states
of finite spectral support.]}
Let $I$ be an instant-class of a lawful branch and consider the
orbits of the licensed unitary family $\{e^{-itH_B}\}$ within the
realization over $I$.  Exactly one of the following holds, and
which one is determined by the declared probe family (whether some
admissible probe fails to commute with the unitary family on
$I$):
\begin{enumerate}[label=\textup{(\Roman*)}]
  \item \emph{(Gauge horn.)}  Every admissible probe is constant
    along unitary orbits over $I$.  Then orbit points lie in a
    single quotient class: the unitary family induces no certified
    change, acts as admissible re-description, and is pure
    presentation.  A reversible clock does not exist because there
    is no licensed motion to measure; unitary ``time'' on $I$ is
    unphysical in the NFC sense.
  \item \emph{(Phase horn.)}  Some admissible probe separates
    orbit points.  Then distinct orbit points realize distinct
    quotient classes, and the orbit motion is certified
    \emph{zero-defect} change within the instant: each induced
    step is probe-witnessed, invertible (the family contains
    $e^{+itH_B}$), and class-count preserving, hence carries no
    ledger event.  For orbit states of finite spectral support
    $\{\lambda_1, \dots, \lambda_r\}$, the certified orbit motion
    is periodic or quasi-periodic with frequency module generated
    by the licensed spectral differences
    $\lambda_i - \lambda_j$, and there exists a \emph{phase
    coordinate} on each orbit --- a toral coordinate of dimension
    the rank of the frequency module, cyclic in the rank-one case
    --- canonical up to phase origin and orientation in each
    generator.  The second clock exists, and it is cyclic/toral,
    not linear: its gauge group is origin~$+$~orientation per
    phase, in contrast to the affine unit~$+$~origin gauge of the
    thermodynamic clock
    (Thm.~\ref{thm:vrp-time-linear}(2)).
\end{enumerate}
\end{theorem}

\begin{proof}
The disjunction is exhaustive and exclusive by inspection of the
declared probe family.  \textbf{Horn (I):} the observational
quotient is generated by admissible probe outcomes; if every
outcome is constant along an orbit, all orbit points are
outcome-indistinguishable, hence lie in one quotient class.  A
quotient-trivial action induces no certified state change, and a
transformation inducing none is an admissible re-description by
the definition of the quotient; there is no licensed change, so
no clock has an object.

\textbf{Horn (II):}  Let $p$ be a separating probe.  Orbit points
with distinct $p$-outcomes realize distinct quotient classes, and
each transition is certified by $p$ itself.  Zero defect: the
induced steps occur at fixed refinement depth (no probe-family
refinement is performed by the motion), are invertible within the
licensed family, and preserve class counts; the defect ledger
records class-splitting and coarsening events only
(Thm.~\ref{thm:vrp-entropy-clock}(1)), and the motion performs
neither, so its ledger increment is zero and the entire orbit
lies within one instant --- consistent with, and required by,
Thm.~\ref{thm:vrp-time-linear}(3).  For the phase coordinate:
with finite spectral support, the spectral theorem gives
\[
  e^{-itH_B}\psi \;=\; \sum_{i=1}^{r} e^{-it\lambda_i}\,
  P_{\lambda_i}\psi,
\]
so every probe outcome along the orbit is a function of the
relative phases $e^{-it(\lambda_i - \lambda_j)}$ --- the orbit's
certified image is a (quasi-)periodic trajectory whose frequency
module is generated by the differences
$\lambda_i - \lambda_j$, which are licensed data
(differences of spectral values of the licensed generator, each a
quotient-visible invariant in the sense of
Thm.~\ref{thm:vrp-mass}'s threshold clause).  The map
$t \mapsto (e^{-it(\lambda_i - \lambda_j)})_{i<j}$ descends to a
coordinate on the closure of the orbit's certified image: a torus
of dimension the rank of the frequency module over
$\mathbb{Z}$, a circle when the rank is one.  Two such
coordinates built from the same licensed frequency data differ by
a choice of origin on each phase and an orientation of each
generator --- the stated gauge --- by the same telescoping
argument as Thm.~\ref{thm:vrp-time-linear}(2), applied per
frequency.  \qed
\end{proof}

\begin{corollary}[\status{C}]\label{cor:vrp-no-global-unitary-time}
\textup{(No Global Linear Unitary Time, on Either Horn.)}
\textup{[Conditional on
Thm.~\ref{thm:vrp-reversible-dichotomy}.]}
On neither horn is $e^{-itH_B}$ with $t$ a global linear licensed
time variable licensed: on horn~(I) because the family induces no
licensed change at all, on horn~(II) because the licensed
parameter is cyclic/toral phase, and promoting phase to a global
linear time would require canonically aligning phase coordinates
across distinct orbits and instants --- which is exactly the
synchronization question that the entropy clock answers only for
defect-bearing change (Thm.~\ref{thm:vrp-entropy-clock}(3)) and
that no licensed datum currently answers for defect-free change.
$e^{-itH_B}$ remains canon as internal functional calculus
(Lem.~\ref{lem:vrp-ham-lyapunov} license note), now with its
physical reading settled horn-by-horn.
\end{corollary}

\begin{proof}
Horn~(I) is immediate.  On horn~(II), a global linear $t$ would
induce a simultaneous canonical choice of origin and orientation
for the phase coordinate of every orbit; by the gauge clause of
Thm.~\ref{thm:vrp-reversible-dichotomy}(II) no licensed datum on
a single orbit makes this choice, and no cross-orbit licensed
comparison exists for defect-free change --- a canonical global
choice would therefore be exactly the hidden selection that the
licensing predicate's declaration clause excludes.  \qed
\end{proof}

\begin{remark}[\status{R}]\label{rmk:vrp-phase-connection}
\textup{(Observation: Phase Alignment Is Connection-Type Data.)}
What horn~(II) leaves open --- a rule canonically aligning phase
coordinates across nearby orbits and instants --- has the formal
shape of a \emph{connection}: locally, alignment data on fibers of
phases; globally, the obstruction to consistent alignment is
holonomy-type.  The corpus's gauge-theoretic machinery (YM
branch) is the natural eventual home of this observation, which
would, if developed, recover gauge structure as the bookkeeping of
reversible-clock synchronization.  Recorded as an observation
only; no force conferred, no development claimed.
\end{remark}

\begin{remark}[\status{C}]\label{ob:vrp-phase-sync}
\textbf{O\_Phase.Sync (Cross-Orbit Phase Synchronization) ---
Conditionally Discharged by the Synchronization Obstruction.}
\textup{[Conditional on Cor.~\ref{cor:vrp-phase-sync-discharge},
\S\,\ref{sec:vrp-sync}.]}
Phase coordinates synchronize across an orbit family iff the phase
holonomy vanishes on every independent cycle of the interface
nerve --- the obstruction in $H^1(N; T^r)$ supported on
$\beta_1(N)$, with vanishing branch data and automatic when
$\beta_1 = 0$.  The expectation that this obligation is
obstruction-bearing (recorded at surfacing) is confirmed at
theorem level: the obstruction is the holonomy of the
reversible-clock connection, and by
Rmk.~\ref{rmk:vrp-sync-gauge} it is supported on exactly the
cycle rank that carries non-abelian gauge structure.
\end{remark}

\subsection{Characterization of the Energy Form: O\_Energy.Form
Resolved at the Ledger Level}
\label{sec:vrp-energy-form}

The last in-toolkit VRP residual asked for universal existence of
the branch energy form.  The L+29 negative result
(Rmk.~\ref{rmk:vrp-univ-method-scope}) showed why no
class-counting construction can deliver it; this subsection
resolves the obligation the way the reversible-clock question
resolved at L+33 --- by characterization rather than one-sided
answer.  Existence at the ledger-functional level is equivalent
to a declaration: a lawful branch carries an energy form exactly
when it declares weighted transport-coercion structure, and on
such branches the form is \emph{constructed} canonically.  The
operator level is then carried by a hypothesis the corpus has
named since L+24.

\begin{definition}[\status{D}]\label{def:vrp-weighted-coercion}
\textup{(Weighted Transport-Coercion Structure.)}
A lawful branch $B$ carries \emph{weighted transport-coercion
structure} when the ``declared structural data'' of transfer
coercivity (Def.~\ref{def:transfer-coercivity}) takes the
quantitative form of \emph{per-step coercion weights}: an
assignment $w(\sigma) \ge 0$ to each admissible transfer step
$\sigma$, declared with:
\begin{enumerate}[label=\textup{(W\arabic*)}]
  \item\label{wc:chain} \emph{(Chain bound.)}  The accumulated
    coercive change along a composable chain is bounded by the sum
    of its step weights --- the quantitative instance of
    Def.~\ref{def:transfer-coercivity};
  \item\label{wc:gap} \emph{(Coercivity gap.)}  There is
    $c_B > 0$ such that every admissible chain from the
    transport-trivial class to a persistent nontrivial class
    accumulates total weight $\ge c_B$;
  \item\label{wc:refine} \emph{(Refinement stability.)}  Weights
    are stable under admissible refinement: a refined step carries
    the weight of the step it refines, and refinement introduces
    no negative-weight steps.
\end{enumerate}
\end{definition}

\begin{theorem}[\status{C}]\label{thm:vrp-energy-form-char}
\textup{(O\_Energy.Form Resolution: Characterization and Canonical
Construction.)}
\textup{[Conditional on: Thm.~\ref{thm:vrp-energy}; the Book~V
lawfulness discipline; and, for clause~(2), the declared weighted
structure of Def.~\ref{def:vrp-weighted-coercion}.]}
\begin{enumerate}[label=\textup{(\arabic*)}]
  \item \emph{(Necessity, by extraction.)}  If a lawful branch
    carries a branch energy form
    (Def.~\ref{def:branch-energy-form}), then its declared data
    certifies weighted transport-coercion structure: the
    increments of $\mathfrak{h}_B$ along admissible steps are
    per-step weights satisfying \ref{wc:chain} (the form's
    transport law), \ref{wc:gap} (clause~\ref{en:pos}: positive
    cost of persistent nontrivial excitation), and
    \ref{wc:refine} (clause~\ref{en:refine}).  An energy form
    \emph{is} a weight declaration.
  \item \emph{(Sufficiency, by canonical construction.)}
    Conversely, given declared weighted structure, the
    \emph{canonical persistence cost}
    \[
      \mathfrak{h}_B^{\min}(x) \;:=\;
      \inf\Bigl\{\, \textstyle\sum_{\sigma \in \gamma} w(\sigma)
      \;:\; \gamma \text{ an admissible transfer chain from the
      transport-trivial class to } x \,\Bigr\}
    \]
    satisfies clauses~\ref{en:zero}--\ref{en:refine} of
    Def.~\ref{def:branch-energy-form}; clause~\ref{en:lsc} holds
    as the usual bridge clause on declared bridge regimes; and
    $\mathfrak{h}_B^{\min}$ is canonical: it is the
    \emph{maximal} functional subordinate to the declared weights
    (any functional vanishing on the trivial class whose
    per-step increments respect the weights is dominated by it),
    determined by the declaration with no further choice.
  \item \emph{(The residual is a pre-existing named hypothesis,
    not new open content.)}  What the construction does not and
    cannot supply is clause~\ref{en:gen} at operator strength: a
    Hilbert realization in which $\mathfrak{h}_B^{\min}$ is a
    densely defined closable quadratic form.  This is precisely
    the realization hypothesis carried, by name, in the bracket of
    Thm.~\ref{thm:vrp-hamiltonian} since its promotion; YM, GR,
    and NS are the certified models in which it holds.  The
    universal-existence obligation therefore resolves with no new
    obligation surfaced: ledger-level existence is characterized
    and constructed; operator-level existence is conditional on
    declared realization, where it always was.
\end{enumerate}
\end{theorem}

\begin{proof}
\textbf{(1)}  Set $w(\sigma) := $ the ledgered increment of
$\mathfrak{h}_B$ along $\sigma$ (nonnegative in the gain
direction by transfer-coercivity; decrements ride the dissipation
ledger and bound the accumulated change).  \ref{wc:chain} is the
form's transport law (Thm.~\ref{thm:vrp-energy} clause~(3))
telescoped along chains.  \ref{wc:gap}: a chain from trivial to a
persistent nontrivial class $x$ accumulates at least
$\mathfrak{h}_B(x) > 0$ by clauses~\ref{en:zero}--\ref{en:pos};
take $c_B$ the infimum of $\mathfrak{h}_B$ over persistent
nontrivial classes, positive on the stabilized finite quotient
(Thm.~\ref{thm:vrp-entropy-univ}(1) supplies finiteness for
lawful branches).  \ref{wc:refine} is clause~\ref{en:refine}.

\textbf{(2)}  \emph{Well-defined and \ref{en:zero}:}  the trivial
class is reachable from itself at cost $0$, so
$\mathfrak{h}_B^{\min}(\mathrm{triv}) = 0$; for persistent
nontrivial $x$, every admissible chain accumulates $\ge c_B$ by
\ref{wc:gap}, so $\mathfrak{h}_B^{\min}(x) \ge c_B > 0$ --- which
is also \ref{en:pos}, with the zero set exactly the
transport-trivial (zero-cost-reachable) classes.
\emph{\ref{en:refine}:}  under admissible refinement the chain
family for a refined class consists of refinements of coarse
chains (paths must respect the finer quotient), so the infimum is
taken over a smaller constraint set with weights preserved
(\ref{wc:refine}): $\mathfrak{h}_B^{\min}$ is non-decreasing
under refinement --- recorded cost is never spuriously destroyed.
\emph{Transport law:}  chain concatenation gives the subordinate
inequality
$\mathfrak{h}_B^{\min}(y) \le \mathfrak{h}_B^{\min}(x) +
w(x \to y)$ --- the no-spurious-gain face of transfer-coercivity,
and the governing pattern (Thm.~\ref{thm:governing}) in its
coercive instance.  \emph{Maximality and canonicity:}  if
$\mathfrak{h}$ vanishes on the trivial class and its per-step
increments respect the weights, then telescoping along any chain
$\gamma$ to $x$ gives
$\mathfrak{h}(x) \le \sum_{\sigma \in \gamma} w(\sigma)$; taking
the infimum over $\gamma$ yields
$\mathfrak{h} \le \mathfrak{h}_B^{\min}$.  The construction uses
the declaration and nothing else.  \emph{\ref{en:lsc}:}  on
declared bridge regimes, lower semicontinuity of a path-infimum
cost under vanishing refinement defect is the same bridge
argument as Thm.~\ref{thm:vrp-time-smooth}, applied to the weight
ledger in place of the clock ledger.

\textbf{(3)}  A path-infimum cost on a quotient is not, as such,
a quadratic form on a Hilbert space; no licensed datum of the
declaration selects a realization, and selecting one is Book~V
realization-functor territory
(\texttt{def:legitimacy-witness}), declared per branch.  The
hypothesis is already named verbatim in the bracket of
Thm.~\ref{thm:vrp-hamiltonian}; nothing new is open.  \qed
\end{proof}

\begin{remark}[\status{R}]\label{rmk:vrp-why-quadratic}
\textup{(The Question Deliberately Not Registered as an
Obligation.)}
``When does the canonical persistence cost admit a closable
quadratic realization?''\ --- in effect, \emph{why is energy
quadratic?} --- is recorded here at observation strength and
deliberately \emph{not} registered as an [O].  Reason: it is a
realization-selection question (which Hilbert realization a
branch declares, Book~V functor territory), not a theorem target
within the recovery program; registering it as an obligation
would misclassify a declaration choice as an unproved claim.  The
certified models answer it instance-wise: in YM the declared
weights are curvature increments and the canonical cost is
realized by the curvature form on its sector
(\texttt{thm:B1-closable}, YM Branch), with GR and NS as the
parallel instances of Rmk.~\ref{rmk:vrp-energy-instances}.  No
force conferred.
\end{remark}

\subsection{Temperature as the Entropy--Energy Exchange Slope}
\label{sec:vrp-temperature}

With the entropy ledger universal (Thm.~\ref{thm:vrp-entropy-univ})
and the energy form characterized
(Thm.~\ref{thm:vrp-energy-form-char}), the Temperature rung ---
anticipated at the gate note of
Rmk.~\ref{rmk:vrp-entropy-ledger-duality} --- has both inputs
canonical.  Temperature is not a primitive ledger but a
\emph{response coefficient}: the marginal rate at which coercive
cost must be spent to buy distinguishability along an admissible
exchange family.  The construction is choice-free given an exchange
family, and the one classical ingredient it must refuse to import
--- the ensemble --- is refused by construction, not by fiat: no
partition function appears because the slope \emph{is} the
definition.

\begin{definition}[\status{D}]\label{def:vrp-exchange-family}
\textup{(Admissible Exchange Family; Exchange-Coherence.)}
Let $B$ carry the universal entropy ledger $S_B$ and a licensed
energy form $\mathfrak{h}_B$ (its ledger reading $E_B$).  An
\emph{admissible exchange family} on a sector $X \subseteq
\Sigma_B$ is a set of admissible steps each of which changes both
ledgers by recorded increments $(\Delta S, \Delta E)$ with
$\Delta E \neq 0$, closed under the branch's admissible
composition.  The family is \emph{exchange-coherent} if the
marginal ratio
\[
  \beta_B \;:=\; \frac{\Delta S_B}{\Delta E_B}
\]
is the same for all steps between any two states of $X$ --- the
exchange analogue of $\Delta$-coherence
(Def.~\ref{def:vrp-coherent-sector}), i.e.\ path-independence of
the entropy bought per unit coercive cost spent.
\end{definition}

\begin{theorem}[\status{C}]\label{thm:vrp-temperature}
\textup{(VRP-Temperature: the Exchange Slope as Licensed Response
Coefficient.)}
\textup{[Conditional on: Thm.~\ref{thm:vrp-entropy-univ},
Thm.~\ref{thm:vrp-energy} (with a licensed energy form on the
sector); and an exchange-coherent admissible family
(Def.~\ref{def:vrp-exchange-family}).]}
On an exchange-coherent sector the marginal slope $\beta_B =
\partial S_B / \partial E_B$ is a licensed physics variable,
realized through the \emph{certified branch-interface coefficient}
disjunct of \texttt{def:variable-licensing} (Book~VII); where
$\beta_B > 0$, the \emph{branch temperature} is
$T_B := \beta_B^{-1}$.  The five recovery clauses:
\begin{enumerate}[label=\textup{(\arabic*)}]
  \item \emph{Probe signature:} the admissible exchange steps and
    the two ledger increments they record.
  \item \emph{Quotient invariant:} $\beta_B$ is a ratio of two
    quotient-visible ledger increments
    (Thms.~\ref{thm:vrp-entropy-univ}, \ref{thm:vrp-energy}),
    constant on the sector by exchange-coherence, hence invariant
    under admissible re-description.
  \item \emph{Transport law:} along lawful interfaces the slope
    transforms by the ratio of the two ledgers' transfer behavior
    --- entropy lax-functorially with no exchange coefficient
    (Thm.~\ref{thm:vrp-entropy-transfer}), energy with its
    declared exchange coefficient
    (Thm.~\ref{thm:vrp-energy-transfer}) --- so $\beta_B$ carries
    exactly the inverse energy exchange coefficient and no
    independent transfer freedom.
  \item \emph{Coherence:} the slope is a state function precisely
    on exchange-coherent sectors; the discharge of the registered
    obligation O\_Temp.Slope is this clause.
  \item \emph{Anti-smuggling audit (O\_Temp.Ensemble, discharged
    by construction):} no canonical ensemble, partition function,
    heat bath, Boltzmann factor, or fitted $\beta$ appears.
    $\beta_B$ is the recorded marginal exchange ratio of two
    licensed ledgers; the ensemble is not refused by assumption
    but is simply never introduced, because the response slope is
    defined directly.
\end{enumerate}
\emph{Sign clause (by design, not omission):} positivity
$\beta_B > 0$ --- coercive cost and distinguishability increasing
together --- is sector data, exactly as gap positivity is branch
data in Thm.~\ref{thm:vrp-mass}.  Exchange families that trade the
ledgers oppositely yield $\beta_B \le 0$ (the population-inversion
regime); the theorem licenses the slope at whatever sign the
recorded increments give, and reserves the name ``temperature''
for the positive branch.
\end{theorem}

\begin{proof}
Quotient-visibility of each increment is
Thm.~\ref{thm:vrp-entropy-univ}(2) for $S_B$ and
Thm.~\ref{thm:vrp-energy} clause~(2) for $E_B$; their ratio is
therefore quotient-determined wherever it is well-defined, and
exchange-coherence (Def.~\ref{def:vrp-exchange-family}) is exactly
the condition that the ratio is a single sector constant ---
established by the same telescoping argument as
Thm.~\ref{thm:vrp-time-linear}(1): a path-dependent marginal ratio
is inconsistent with any state function of the two ledgers, and on
a coherent sector the common ratio defines $\beta_B$.  The
transport law is the quotient of the two transfer theorems: under
a lawful interface $\Delta S_B$ is carried with no exchange
coefficient (Thm.~\ref{thm:vrp-entropy-transfer}) and $\Delta E_B$
with coefficient $g_I$ (Thm.~\ref{thm:vrp-energy-transfer}), so
$\beta_B \mapsto g_I^{-1}\beta_B$ --- the certified
branch-interface coefficient disjunct of the licensing predicate
is met with $\beta_B$'s sole transfer freedom being the inverse
energy coefficient.  The declaration clause is satisfied by the
itemization in the statement, including the explicit sign clause
and the anti-smuggling audit.  No ensemble construct is used at
any step. \qed
\end{proof}

\begin{remark}[\status{R}]\label{rmk:vrp-temperature-dissipation}
\textup{(Instance Pattern: the Dissipation Ledger Reads as a
Temperature.)}
The L+27 dissipation ledger
(Lem.~\ref{lem:vrp-ham-lyapunov}) is the natural exchange family:
along $H_B$-compatible continuation the coercive cost
$\mathfrak{h}_B$ decreases by ledgered amounts, and where that
spent cost registers as an entropy increase the recorded ratio
$\Delta S_B / \Delta E_B$ is precisely $\beta_B$ --- a
fluctuation--dissipation-shaped reading in which dissipated
coercive cost and produced distinguishability are the two sides of
one ledgered step.  In the NS safe-tail regime the contraction
$\mathfrak{D}_{n+1} \le \kappa\,\mathfrak{D}_n$
(\texttt{cor:NS61-discharged}, NS Branch) supplies a sign-definite
exchange family of exactly this shape.  Recorded as an instance
pattern; no force conferred and no equilibrium claimed.
\end{remark}

\begin{remark}[\status{C}]\label{ob:vrp-temp-equil}
\textbf{O\_Temp.Equil (the Zeroth Law) --- Conditionally
Discharged by the Synchronization Obstruction.}
\textup{[Conditional on Cor.~\ref{cor:vrp-temp-equil-discharge},
\S\,\ref{sec:vrp-sync}.]}
The zeroth law holds as transitivity of thermal equilibrium along
composable interfaces (by exchange-coefficient multiplicativity),
and global temperature synchronization across a family is
characterized: it holds iff the exchange-coefficient holonomy
vanishes on every independent cycle of the interface nerve, the
obstruction in $H^1(N; \mathbb{R}_{>0})$ supported on
$\beta_1(N)$, automatic when $\beta_1 = 0$.  Whether the
obstruction vanishes for a given branch's nerve is branch data,
exactly as gap positivity is branch data
(Thm.~\ref{thm:vrp-mass}).  The shell's O\_Temp.Kinetic
(kinetic-theory reading) remains out of toolkit and is not
registered as an in-corpus [O].
\end{remark}

\subsection{Charge as Licensed Conserved Quotient Invariant}
\label{sec:vrp-charge}

The conserved charge space $\mathcal{Q}_\lambda = \ker(T_\partial -
\lambda I)^{\!\top}$ (Def.~\ref{def:conserved-charges}) and its
compatibility gate (Prop.~\ref{prop:charge-compatibility-gate},
unconditional) are already canonical.  This subsection licenses the
charge as a recovered physics variable, binding it to the continuation
order (conservation) and to the type-transition operator
(quantization), with the gate supplying the transfer law at no extra
cost.  Charge is the most directly NFC-native of the recovered
variables after entropy: it is a left-kernel of an operator the canon
already carries, not a new construction.

\begin{definition}[\status{D}]\label{def:vrp-charge-observable}
\textup{(Charge Observable.)}
For $\psi \in \mathcal{Q}_\lambda$ (Def.~\ref{def:conserved-charges}),
the \emph{charge observable} is the linear functional
\[
  q_\psi : \mathcal{O}_B \to \mathbb{R}, \qquad
  q_\psi(x) := \langle \psi, x \rangle,
\]
the left-eigenvector pairing of the observable state $x$ with the
conserved direction $\psi$.  The \emph{charge spectrum} is the set
$\{\lambda\}$ of eigenvalues of $T_\partial$ carrying a nonzero
$\mathcal{Q}_\lambda$.
\end{definition}

\begin{theorem}[\status{C}]\label{thm:vrp-charge}
\textup{(VRP-Charge: Conserved Charge as Licensed Quotient Invariant.)}
\textup{[Conditional on: the Book~II type-transition structure
(Def.~\ref{def:T-partial}) and Def.~\ref{def:conserved-charges}.]}
The charge observable $q_\psi$ (Def.~\ref{def:vrp-charge-observable})
is a licensed physics variable in the sense of
\texttt{def:variable-licensing} (Book~VII), realized through the
quotient-visible-invariant disjunct.  The five recovery clauses:
\begin{enumerate}[label=\textup{(\arabic*)}]
  \item \emph{Probe signature:} the left-eigenvector pairing
    $\langle\psi,\cdot\rangle$, evaluated by the declared admissible
    probes that read the type-transition structure.
  \item \emph{Quotient invariant:} $q_\psi$ is conserved under every
    admissible transfer map (Def.~\ref{def:conserved-charges}), so its
    value is constant on observational-equivalence classes and
    invariant under admissible re-description.
  \item \emph{Transport law:} supplied unconditionally by the
    compatibility gate (Prop.~\ref{prop:charge-compatibility-gate}):
    $q_\psi$ is the exact obstruction to interface solvability
    ($\mathcal{A}_\lambda u = J_w$ solvable iff $q_\psi(J_w)=0$ for all
    $\psi\in\mathcal{Q}_\lambda$), so charge transfers across lawful
    interfaces by the Fredholm pairing with no independent transfer
    freedom.
  \item \emph{Additivity:} $q_\psi$ is linear, hence additive over
    independent (direct-sum) subsystems --- the ledger property of a
    physical charge.
  \item \emph{Anti-smuggling audit:} no Noether theorem, Lagrangian,
    action principle, continuous symmetry group, or gauge potential is
    primitive; charge is the left-kernel of the licensed
    type-transition operator, recovered from canonical structure, not
    assumed from a symmetry postulate.
\end{enumerate}
\end{theorem}

\begin{proof}
Conservation under admissible transfer is Def.~\ref{def:conserved-charges}
(left eigenvectors of $T_\partial$ are fixed, up to the eigenvalue
scaling, under the transfer maps generated by $T_\partial$), giving
clause~(2): the pairing is constant on equivalence classes and
re-description-invariant.  Clause~(3) is
Prop.~\ref{prop:charge-compatibility-gate} verbatim --- an
unconditional result, so the transport law is inherited at full
force.  Linearity of $\langle\psi,\cdot\rangle$ gives clause~(4).
Clauses~(1) and~(5) are the declaration: the probe signature is the
pairing, and the anti-smuggling audit holds because every ingredient
($T_\partial$, its left kernel, the pairing) is canonical NFC
structure with no symmetry primitive introduced.  The
quotient-visible-invariant disjunct of the licensing predicate is thus
satisfied. \qed
\end{proof}

\begin{corollary}[\status{C}]\label{cor:vrp-charge-conservation}
\textup{(Charge Is Conserved Along the Licensed Continuation Order.)}
\textup{[Conditional on Thm.~\ref{thm:vrp-charge},
Thm.~\ref{thm:vrp-time}.]}
For every $\psi\in\mathcal{Q}_\lambda$, the charge $q_\psi$ is constant
along the licensed continuation order $\preceq_t$: if $x \preceq_t y$
then $q_\psi(x) = q_\psi(y)$.  Charge is thereby bound to branch time
as its conserved companion --- a further inter-variable binding (charge
$\leftrightarrow$ order), alongside the generator--order binding (L+27)
and the entropy--order binding (L+28).
\end{corollary}

\begin{proof}
Each generating step of $\preceq_t$ (Def.~\ref{def:branch-time-order})
is an admissible quotient-visible state change, hence realized by an
admissible transfer map, under which $\mathcal{Q}_\lambda$ is conserved
(Def.~\ref{def:conserved-charges}).  By clause~(2) of
Thm.~\ref{thm:vrp-charge} the pairing is unchanged across each step;
induction on chain length gives constancy along $\preceq_t$. \qed
\end{proof}

\begin{theorem}[\status{C}]\label{thm:vrp-charge-quantization}
\textup{(Charge Discreteness in the Finite Regime; Continuum Survival
Open.)}
\textup{[Conditional on Thm.~\ref{thm:vrp-charge}; finite-dimensional
type space $V$ at each refinement stage.]}
On the finite NFC regime the charge spectrum $\{\lambda\}$ is
\emph{discrete}: $T_\partial$ acts on a finite-dimensional type space,
so its spectrum is finite and the admissible charge values are
isolated.  Charge quantization is therefore not an external postulate
but a consequence of the finite type-transition structure --- charges
are discrete because types are.  What is \emph{not} established here is
survival of discreteness under the Book~VI continuum bridge: whether
the isolated charge spectrum persists, or accumulates, in the continuum
encoding is the open obligation O\_Charge.Quant-cont
(Rmk.~\ref{ob:vrp-charge-quant}).
\end{theorem}

\begin{proof}
A linear operator on a finite-dimensional space has finite spectrum;
the eigenvalues carrying nonzero left-kernel $\mathcal{Q}_\lambda$ are
among them, hence finite and isolated.  Discreteness of the admissible
charge values follows immediately.  The continuum clause is explicitly
deferred and proved nothing here. \qed
\end{proof}

\begin{remark}[\status{R}]\label{rmk:vrp-charge-gauge}
\textup{(Charge, the Gauge Template, and Cycle Rank.)}
The conserved charges organize the canonical gauge content: the YM
gauge template $\mathfrak{su}(3)\oplus\mathfrak{su}(2)\oplus
\mathfrak{u}(1)$ (\texttt{prop:YM-status}) and the SM matter quantum
numbers are charge data in the present sense.  By the synchronization
obstruction (Thm.~\ref{thm:vrp-sync-obstruction}) the global
organization of charges across an interface family is a class in
$H^1(N;G)$ supported on $\beta_1(N)$, the same cycle rank that
\texttt{thm:beta1-controls-nonabelian} (\S\,\ref{sec:beta1}) assigns to
non-abelian gauge structure.  Charge, gauge structure, and
synchronization holonomy thus share one support --- the conserved
charge is the local datum whose global gluing the cycle rank governs.
Recorded as the program's reading; no force conferred beyond the cited
theorems.
\end{remark}

\begin{remark}[\status{O}]\label{ob:vrp-charge-quant}
\textbf{Open obligation O\_Charge.Quant-cont (Continuum Survival of
Charge Discreteness).}
Establish whether the discrete charge spectrum of the finite regime
(Thm.~\ref{thm:vrp-charge-quantization}) survives the Book~VI continuum
bridge --- i.e.\ whether admissible charge values remain isolated, or
can accumulate, in the continuum encoding.  This is the charge-rung
instance of the standing continuum-encoding question (cf.\ the energy
LSC clause and the time smooth-encoding clause); it is the
continuum-bridge question for finite-regime spectral discreteness
generally --- it covers both the charge spectrum
(Thm.~\ref{thm:vrp-charge-quantization}) and the frequency module
(Thm.~\ref{thm:vrp-frequency}, Rmk.~\ref{rmk:vrp-frequency-scope}),
which present the identical deferral --- the finite-regime discreteness
being unconditional within its declared scope in both cases.
\end{remark}

\subsection{Frequency and Phase: Licensing the Spectral-Gap Module
and Its Conjugate}
\label{sec:vrp-frequency}

The reversible-clock dichotomy (Thm.~\ref{thm:vrp-reversible-dichotomy},
phase horn) already constructed, structurally, the phase coordinate on
a unitary orbit and the frequency module generated by the licensed
spectral gaps $\lambda_i - \lambda_j$.  This subsection licenses
\emph{frequency} as a physics variable in its own right --- the
spectral-gap module of the licensed generator --- establishes the
phase--frequency duality (new: the phase coordinate of L+33 was built
but never paired with frequency as conjugate), and binds frequency to
the Hamiltonian spectrum, of which mass (Thm.~\ref{thm:vrp-mass}) is the
bottom gap.  No dimensional constant and no energy--frequency physical
law is introduced; the content is the dimensionless module structure
and its conjugacy.

\begin{definition}[\status{D}]\label{def:vrp-frequency-module}
\textup{(Frequency Module.)}
On an orbit/sector $I$ of finite spectral support
$\{\lambda_1,\dots,\lambda_r\}$ of the licensed generator $H_B$
(Thm.~\ref{thm:vrp-hamiltonian}), the \emph{frequency module} is the
additive subgroup
\[
  \Omega_B(I) \;:=\; \bigl\langle\, \lambda_i - \lambda_j \;:\;
  1 \le i,j \le r \,\bigr\rangle \;\subseteq\; \mathbb{R},
\]
generated by the spectral gaps.  Its rank is the toral dimension of the
phase coordinate of Thm.~\ref{thm:vrp-reversible-dichotomy}(II).
\end{definition}

\begin{theorem}[\status{C}]\label{thm:vrp-frequency}
\textup{(VRP-Frequency: the Spectral-Gap Module as Licensed Variable.)}
\textup{[Conditional on: Thm.~\ref{thm:vrp-hamiltonian},
Thm.~\ref{thm:vrp-mass} (for the spectral-invariance of gaps); finite
spectral support on the sector.]}
The frequency module $\Omega_B(I)$ (Def.~\ref{def:vrp-frequency-module})
is a licensed physics variable through the quotient-visible-invariant
disjunct of \texttt{def:variable-licensing} (Book~VII).  The five
recovery clauses:
\begin{enumerate}[label=\textup{(\arabic*)}]
  \item \emph{Probe signature:} the spectral gaps resolved by the
    admissible probes on the sector (the same probes that, on the phase
    horn, separate orbit points,
    Thm.~\ref{thm:vrp-reversible-dichotomy}(II)).
  \item \emph{Quotient invariant:} the gaps $\lambda_i-\lambda_j$ are
    differences of spectral values of the licensed generator, each a
    quotient-visible invariant (the spectral-threshold argument of
    Thm.~\ref{thm:vrp-mass} applied to every level, not only the
    bottom), so $\Omega_B(I)$ is fixed under admissible re-description.
  \item \emph{Transport law:} the generator transfers by uniqueness
    (Prop.~\ref{prop:vrp-transfer-reduction}), carrying its spectrum
    and hence the gap module; frequency transfers with the Hamiltonian
    and inherits no independent transfer freedom.
  \item \emph{Module structure (additivity):} composing admissible
    transitions adds their gaps, so $\Omega_B(I)$ is closed under
    addition --- the ledger property, here a $\mathbb{Z}$-module of
    frequencies.
  \item \emph{Anti-smuggling audit:} no Planck constant, no
    $E=\hbar\omega$, no oscillator, wave, or periodic-motion primitive,
    and no external clock; frequency is the spectral-gap structure of
    the licensed generator, recovered, not posited.
\end{enumerate}
\end{theorem}

\begin{proof}
Clause~(2) is the per-level spectral-invariance: each $\lambda_i$ is a
spectral value of the licensed self-adjoint $H_B$ on a quotient-visible
sector, invariant under admissible re-description by the same argument
that makes the bottom gap $m^2$ invariant in Thm.~\ref{thm:vrp-mass};
differences and the group they generate are then invariant.
Clause~(3): $H_B$ is the unique generator of the transferred closed
form (Prop.~\ref{prop:vrp-transfer-reduction}), so its spectrum, gaps,
and gap module are carried by the same uniqueness.  Clause~(4) is
closure of $\langle\lambda_i-\lambda_j\rangle$ under addition.
Clauses~(1),(5) are the declaration, the audit holding because every
ingredient is spectral data of the licensed generator with no
dimensional or dynamical primitive introduced.  The
quotient-visible-invariant disjunct is satisfied. \qed
\end{proof}

\begin{theorem}[\status{C}]\label{thm:vrp-phase-frequency-duality}
\textup{(Phase--Frequency Duality.)}
\textup{[Conditional on Thm.~\ref{thm:vrp-frequency},
Thm.~\ref{thm:vrp-reversible-dichotomy}(II).]}
On a phase-horn sector the phase coordinate is a torus $T^r$
(Thm.~\ref{thm:vrp-reversible-dichotomy}(II)) and the frequency module
$\Omega_B(I)$ is its character lattice: $\Omega_B(I) \cong
\widehat{T^r} \cong \mathbb{Z}^r$, and phase and frequency are
\emph{conjugate} in the Pontryagin sense --- each admissible probe
outcome along the orbit is a finite combination of characters
$e^{i\langle \omega, \theta\rangle}$, $\omega\in\Omega_B(I)$,
$\theta\in T^r$.  The duality pairing is the licensed,
re-description-invariant pairing of the two licensed variables.
\emph{Scope clause:} this is abelian Pontryagin duality (torus
$\leftrightarrow$ character lattice), \emph{not} a quantum uncertainty
relation and \emph{not} an energy--frequency identification; no
$\hbar$, no $\Delta\theta\,\Delta\omega$ bound, and no dimensional
constant is asserted or licensed.
\end{theorem}

\begin{proof}
The phase coordinate descends to $T^r$ with $r = \operatorname{rank}
\Omega_B(I)$ (Thm.~\ref{thm:vrp-reversible-dichotomy}(II)); the
characters of $T^r$ are indexed by $\mathbb{Z}^r$, which is exactly the
frequency module under the gap generators (Def.~\ref{def:vrp-frequency-module}).
Each orbit observable is, by the spectral expansion
$e^{-itH_B}\psi=\sum_i e^{-it\lambda_i}P_{\lambda_i}\psi$, a finite
combination of the relative-phase characters $e^{-it(\lambda_i-\lambda_j)}$,
i.e.\ a function on $T^r$ with frequencies in $\Omega_B(I)$.  Pontryagin
duality of $T^r$ and $\mathbb{Z}^r$ gives the conjugacy; both sides are
licensed (phase by L+33, frequency by Thm.~\ref{thm:vrp-frequency}) and
the pairing is re-description-invariant.  The scope clause is a
restriction on the claim, introducing nothing. \qed
\end{proof}

\begin{corollary}[\status{R}]\label{cor:vrp-frequency-energy-binding}
\textup{(Frequency and Mass Are Two Readings of One Spectrum.)}
Frequency is bound to the Hamiltonian as its gap structure: $\Omega_B$
is generated by differences of the spectrum whose bottom gap on the
nontrivial sector is the mass $m^2$ (Thm.~\ref{thm:vrp-mass}).  Mass is
the bottom gap; the frequency module is the full gap lattice; both are
quotient-visible spectral invariants of the one licensed generator.
This is the fifth inter-variable binding (frequency $\leftrightarrow$
generator/mass), and it records that the recovered spectral variables
--- mass and frequency --- are not independent posits but two readings
of $\operatorname{Spec}(H_B)$.  No force beyond the cited theorems.
\end{corollary}

\begin{remark}[\status{R}]\label{rmk:vrp-frequency-scope}
\textup{(Synchronization, Continuum Survival, and What Stays Out of
Toolkit.)}
Three honest scope notes.  \emph{(i) Cross-orbit comparison} of phase
and frequency across an interface family is the synchronization
obstruction already characterized (Thm.~\ref{thm:vrp-sync-obstruction},
Cor.~\ref{cor:vrp-phase-sync-discharge}): a class in $H^1(N;T^r)$ on
$\beta_1$; no new obligation.  \emph{(ii) Continuum survival} of the
discrete frequency module is the frequency-rung instance of the same
continuum-bridge question registered for charge
(Rmk.~\ref{ob:vrp-charge-quant}, O\_Charge.Quant-cont): on the finite
regime $\Omega_B$ is discrete because spectral support is finite;
whether discreteness survives the Book~VI bridge is that one shared
question, not a separate obligation.  \emph{(iii) Out of toolkit:} the
physical energy--frequency law $E=\hbar\omega$ requires a dimensional
quantization constant $\hbar$ that no licensed variable supplies; the
in-toolkit content is the dimensionless module and its duality, and the
$\hbar$-identification is explicitly outside the recovery program.
\end{remark}

%% ---------------------------------------------------------------
\section{Synchronization of Recovered Variables: One Obstruction,
Supported on Cycle Rank}
\label{sec:vrp-sync}
%% ---------------------------------------------------------------

The recovery program leaves a single in-corpus frontier: two
obligations --- cross-orbit phase synchronization
(\texttt{ob:vrp-phase-sync}, from the reversible-clock dichotomy)
and the temperature zeroth law (\texttt{ob:vrp-temp-equil}) ---
that the surfacing passes independently flagged as
connection-shaped.  This section proves they are one obstruction.
Each asks when a locally canonical quantity, defined up to a
structure-group action and related across lawful interfaces by a
multiplicative coefficient, glues into a single global quantity;
the answer is cohomological, the obstruction lives in the first
cohomology of the interface nerve, and --- by the canon's own
cycle-rank theorem (\S\,\ref{sec:beta1}) --- it is supported
exactly on the complexity parameter $\beta_1$ that already governs
non-abelianity.  Synchronization is free precisely where the gauge
structure is trivial.

\begin{definition}[\status{D}]\label{def:vrp-sync-datum}
\textup{(Synchronization Datum over an Interface Nerve.)}
Let $\{B_i\}$ be lawful sectors with lawful interfaces
$I_{ij} : B_i \to B_j$ on overlaps, forming the
\emph{interface nerve} $N$ (vertices the sectors, edges the
lawful interfaces, triangles the certified triple overlaps).
A \emph{synchronization datum} valued in an abelian structure
group $G$ assigns:
\begin{enumerate}[label=\textup{(\roman*)}]
  \item to each sector a locally canonical invariant determined up
    to the action of $G$ (its \emph{gauge}: origin-and-orientation
    for a phase, multiplicative scale for an exchange slope);
  \item to each interface a transition element
    $g_{ij} \in G$ relating the two local invariants, with
    $g_{ji} = g_{ij}^{-1}$;
\end{enumerate}
such that the transitions are \emph{consistent on triple
overlaps}: $g_{ij}\, g_{jk}\, g_{ki} = 1$ on every triangle of
$N$ (the \emph{cocycle condition}).  A datum is
\emph{synchronizable} when there exist gauge choices
$a_i \in G$ with $g_{ij} = a_i\, a_j^{-1}$ for all edges --- i.e.\
the local invariants glue, in a common gauge, into one global
quantity.
\end{definition}

\begin{theorem}[\status{C}]\label{thm:vrp-sync-obstruction}
\textup{(The Synchronization Obstruction is a Cohomology Class
Supported on $\beta_1$.)}
\textup{[Conditional on: a synchronization datum
(Def.~\ref{def:vrp-sync-datum}) over a finite interface nerve $N$
with abelian structure group $G$.]}
\begin{enumerate}[label=\textup{(\arabic*)}]
  \item \emph{(Cohomological characterization.)}  The datum is
    synchronizable if and only if its transition family is a
    coboundary; equivalently, its class
    $[g] \in H^1(N; G)$ vanishes.  This class is the complete
    obstruction: synchronization holds iff $[g] = 0$.
  \item \emph{(Support on cycle rank.)}  For connected $N$,
    $H^1(N; G) \cong \mathrm{Hom}(H_1(N), G)$, and $H_1(N)$ is
    free of rank $\beta_1(N)$.  Hence the obstruction is carried
    by $\beta_1(N)$ independent generators --- one holonomy per
    independent cycle of the interface nerve --- and vanishes
    identically when $\beta_1(N) = 0$ (tree-like interface
    families): on a simply-connected nerve every consistent local
    datum synchronizes, with the common gauge fixed by transport
    along the spanning tree, uniquely up to one global $G$-choice.
  \item \emph{(Holonomy form.)}  For each cycle
    $\gamma = (i_0 \to i_1 \to \cdots \to i_r \to i_0)$ of $N$ the
    obstruction is the holonomy
    $\mathrm{hol}(\gamma) = g_{i_0 i_1} g_{i_1 i_2} \cdots
    g_{i_r i_0} \in G$; the datum synchronizes iff
    $\mathrm{hol}(\gamma) = 1$ on a generating set of cycles.
\end{enumerate}
\end{theorem}

\begin{proof}
\textbf{(1)}  If $g_{ij} = a_i a_j^{-1}$ then on any triangle the
product telescopes to $1$, so a coboundary is a cocycle (the datum
is consistent) and is synchronizable by the gauges $a_i$.
Conversely, synchronizability is by definition the existence of
such $a_i$, i.e.\ that the cocycle $g$ is a coboundary; the
quotient of cocycles by coboundaries is $H^1(N;G)$ by definition,
so synchronizability is exactly the vanishing of $[g]$.

\textbf{(2)}  For a finite connected $1$-complex (and, on triple
overlaps, $2$-complex) with abelian $G$, the universal-coefficient
computation gives $H^1(N;G) \cong \mathrm{Hom}(H_1(N), G)$;
$H_1(N)$ of a finite complex is free abelian of rank the first
Betti number $\beta_1(N)$.  A homomorphism out of a free group of
rank $\beta_1$ is determined by its values on $\beta_1$
generators, so the obstruction has exactly $\beta_1(N)$
independent components.  When $\beta_1 = 0$ the group $H_1$ is
trivial, $H^1 = 0$, and every cocycle is a coboundary: fix any
spanning tree, set $a_{i_0} = 1$ at a root, and define $a_i$ by
the transition product along the unique tree path; consistency
(the cocycle condition) makes this independent of how non-tree
edges would route, so $g_{ij} = a_i a_j^{-1}$ on all edges.

\textbf{(3)}  The holonomy around $\gamma$ is the coboundary
defect of the cycle; by~(1)--(2) the class $[g]$ is determined by
the holonomies on a generating set of $H_1(N)$, and
$[g] = 0$ iff all such holonomies are trivial. \qed
\end{proof}

\begin{corollary}[\status{C}]\label{cor:vrp-temp-equil-discharge}
\textup{(Zeroth Law and Temperature Synchronization.)}
\textup{[Conditional on Thm.~\ref{thm:vrp-sync-obstruction};
Thm.~\ref{thm:vrp-temperature}; and exchange-coefficient
multiplicativity, Thm.~\ref{thm:vrp-energy-transfer}(3).]}
Take $G = (\mathbb{R}_{>0}, \times)$ and transition elements the
energy exchange coefficients $g_{ij} = g_{I_{ij}}$ relating the
slopes $\beta_{B_i}$ (Thm.~\ref{thm:vrp-temperature}(3)).
Multiplicativity makes this a synchronization datum.  Then:
\begin{enumerate}[label=\textup{(\arabic*)}]
  \item \emph{(Zeroth law, transitivity.)}  Along any path of
    composable interfaces, equality of $g$-corrected slopes is
    transitive --- thermal equilibrium is an equivalence relation
    on path-connected sectors --- by multiplicativity composing
    the corrections.
  \item \emph{(Global common temperature.)}  A family of sectors
    admits a single common temperature (one global slope in a
    common gauge) iff the exchange-coefficient holonomy vanishes
    on every independent cycle of the interface nerve, i.e.\ iff
    the obstruction in $H^1(N; \mathbb{R}_{>0})$ vanishes;
    automatic when $\beta_1(N) = 0$.
\end{enumerate}
This conditionally discharges O\_Temp.Equil: the zeroth law holds
as transitivity, and global synchronization is characterized with
the obstruction named; whether the obstruction vanishes for a
given branch's interface nerve is branch data, exactly as gap
positivity is branch data in Thm.~\ref{thm:vrp-mass}.
\end{corollary}

\begin{proof}
Multiplicativity (Thm.~\ref{thm:vrp-energy-transfer}(3)) is the
cocycle condition for $g_{ij} = g_{I_{ij}}$, so
Def.~\ref{def:vrp-sync-datum} is met.  (1) is the path-level
specialization (composites multiply); (2) is
Thm.~\ref{thm:vrp-sync-obstruction}(1)--(2) for
$G = \mathbb{R}_{>0}$. \qed
\end{proof}

\begin{corollary}[\status{C}]\label{cor:vrp-phase-sync-discharge}
\textup{(Phase Synchronization and the Reversible-Clock
Holonomy.)}
\textup{[Conditional on Thm.~\ref{thm:vrp-sync-obstruction};
Thm.~\ref{thm:vrp-reversible-dichotomy} on the phase horn; and
local consistency of phase alignment --- the spectral-difference
data $\lambda_i - \lambda_j$ agreeing on overlaps, which is the
cocycle condition for the phase transitions.]}
Take $G$ the phase torus $T^r$ (the structure group of the phase
coordinate, $r$ the frequency-module rank) and transition elements
the relative phase alignments on overlapping orbits.  Then phase
coordinates synchronize across the family iff the phase holonomy
vanishes on every independent cycle of the interface nerve ---
the obstruction in $H^1(N; T^r)$, supported on $\beta_1(N)$.  This
conditionally discharges O\_Phase.Sync: synchronization is
characterized with the holonomy named; its vanishing is branch
data.  In particular the reversible clock globalizes on tree-like
orbit families ($\beta_1 = 0$) and is obstructed exactly by the
cycle structure of the interface nerve.
\end{corollary}

\begin{proof}
On the phase horn the per-orbit coordinate is canonical up to
origin-and-orientation (Thm.~\ref{thm:vrp-reversible-dichotomy}(II)),
which is the $G = T^r$ gauge; relative alignments on overlaps are
the transition elements, a cocycle by the assumed local
consistency.  Apply Thm.~\ref{thm:vrp-sync-obstruction}. \qed
\end{proof}

\begin{remark}[\status{R}]\label{rmk:vrp-sync-gauge}
\textup{(The Obstruction Vanishes Exactly Where the Gauge
Structure Is Trivial.)}
Thm.~\ref{thm:vrp-sync-obstruction}(2) places the synchronization
obstruction on $\beta_1(N)$.  The canon's cycle-rank theorem
(\texttt{thm:beta1-controls-nonabelian}, \S\,\ref{sec:beta1})
places non-abelianity on the same $\beta_1$:
$\beta_1 = 0$ is the abelian sector, and $\beta_1$ is the rank of
the non-abelian structure.  These are the same support.  The
reading anticipated at the reversible-clock pass
(\texttt{rmk:vrp-phase-connection}) is thereby a theorem-level
identification: the failure of recovered variables to synchronize
is the holonomy of a connection on the interface nerve, the
connection is flat (consistent) by the cocycle condition, its
holonomy lives on the cycles, and gauge structure is precisely the
bookkeeping of that holonomy.  Synchronization is free exactly
when there is no gauge structure to obstruct it.  Recorded as the
program's reading of its own frontier; no force conferred beyond
the cited theorems.
\end{remark}

\section{$\beta_1$ Cycle Rank as Canonical Complexity Parameter}
\label{sec:beta1}
%% ---------------------------------------------------------------

\begin{definition}[\status{D}]\label{def:cycle-rank}
\textup{(Cycle Rank $\beta_1$ of the LS-2 Adjacency Graph.)}
Let $\Gamma_{\mathrm{LS2}}$ be the LS-2 adjacency graph on the
stabilized collar type alphabet $\Sigma_\partial$ (Book~II):
the graph whose vertices are collar type classes and whose edges
are admissible single-step collar transitions.  The
\emph{cycle rank} (first Betti number) is
\[
  \beta_1(\Gamma_{\mathrm{LS2}})
  \;:=\; |\mathrm{Edges}| - |\mathrm{Vertices}|
  + \text{(connected components)}.
\]
\end{definition}

\begin{theorem}[\status{C}]\label{thm:beta1-controls-nonabelian}
\textup{($\beta_1$ Controls Non-Abelianity.)}
\textup{[Conditional on $K_0=7$, LS-2, Phase-A.]}
The cycle rank $\beta_1$ of the LS-2 adjacency graph determines
the complexity of the non-abelian collar-local algebra:
\begin{enumerate}[label=\textup{(\arabic*)}]
  \item $\dim \mathfrak{g}' \leq \beta_1$ (Loop Rank bound);
  \item $\beta_1 = 0$: abelian sector only;
  \item $\beta_1 = 1$: non-abelian sector contains $A_1
    = \mathfrak{su}(2)$ ($|\Phi| = 2$);
  \item $\beta_1 = 2$: non-abelian sector contains $A_2
    = \mathfrak{su}(3)$ ($|\Phi| = 6$).
\end{enumerate}
At $K_0=7$: the color sector ($B_{+1}$ block) has $\beta_1 = 2$;
the weak sector (reduced active) has $\beta_1 = 1$; the
background ($B_0$) has $\beta_1 = 0$.
\end{theorem}

\begin{proof}
The dimension bound $\dim \mathfrak{g}' \leq \beta_1$ follows
from the Collar-Local Lie Theory chain (SM Branch,
Thm.~\ref{thm:sm-collar-lie-chain}, Step~5: Loop Rank).
The classification by $\beta_1$ is the rank classification of
compact semisimple Lie algebras: rank $= \beta_1$ selects
the unique compact simple type with $|\Phi| \leq \rho = 8$.
For $K_0=7$: the $d_{S_v}$-split gives three blocks with cycle
ranks $2, 1, 0$ respectively (from the NF2 class space structure,
Book~II, Def.~\ref{def:NF2-blocks}).
\qed
\end{proof}

\begin{remark}[\status{R}]
$\beta_1$ serves as a canonical complexity measure throughout the
NFC program:
\begin{itemize}
  \item In the YM branch: $\beta_1 = 2$ for the color sector
    drives the $A_2$ gauge algebra;
  \item In the NS branch: $\beta_1$ of the NS obstacle adjacency
    graph measures the difficulty of the regularity problem
    (higher $\beta_1$ = more independent obstruction modes);
  \item In the SM super-branch: combined $\beta_1^{\mathrm{total}}
    = 3$ saturates the $\rho = 8$ obstruction bound
    ($|\Phi(A_2)| + |\Phi(A_1)| = 6 + 2 = 8$).
\end{itemize}
\end{remark}

\section{Boundary of Book III}
%% ---------------------------------------------------------------

\textbf{What Book~III proves.}
\begin{enumerate}
  \item Persistent observables are well-defined, quotient-visible,
        and transfer-compatible.
  \item Admissible transfer maps compose, are witness-preserving,
        and respect the observational quotient.
  \item The interface-relative stability functional satisfies the
        Unified Coercive-Transport Inequality at every admissible
        enlargement step.
  \item Finite balance data compose along composable transfer
        chains.
  \item Under extra coercivity hypotheses, observational drift is
        geometrically controlled.
  \item Observational-transfer objects form a well-defined
        comparison arena.
  \item Depth-sum convergence for Van Hove exhaustions
        (Thm~\ref{thm:depth-sum}): weighted boundary contributions
        vanish in the thermodynamic limit.
  \item The ICDC Schema is recorded as an architectural scholium
        (Sch.~\ref{sch:icdc-abstract}): the abstract UCTI
        scale-limit formalization.
  \item Two topological conservation scholia
        (Schs.~\ref{sch:topological-conservation}--\ref{sch:hopf-helicity}):
        discrete-time $Q_{\mathrm{simp}}$ conservation and Hopf
        helicity conservation under admissible kernel evolution.
\end{enumerate}

\textbf{What Book~III does not prove, and may not be assumed here.}
\begin{enumerate}
  \item \emph{Universal source structure.}  The existence of a
        universal source object $U \in \mathrm{POT}$ is a theorem of
        Book~IV.  Book~III builds the transport machinery that makes
        source completion meaningful; it does not prove source
        completion itself.
  \item \emph{Branch specializations.}  The Yang--Mills
        specialization ($\Theta_n = 1$, gap-subcritical closure),
        the Navier--Stokes specialization ($\Theta_n = \theta < 1$,
        quasi-decay), and the gravity specialization all belong in
        derived branch books.  No specialization of the Unified
        Coercive-Transport Inequality to a specific endpoint is
        canonical in this book.
  \item \emph{Continuum structure.}  No continuum, differential, or
        smooth language is licensed at this stage.  The transport
        inequality is a finite discrete bound.
  \item \emph{Branch closure.}  Transport coercivity at this level is
        part of source-preparation discipline.  Whether it extends to
        closure of any specific branch requires branch-level
        verification.
  \item \emph{Physical endpoint identification.}  No identification
        with gauge fields, fluid dynamics, or geometric structure is
        claimed.
\end{enumerate}

%% ---------------------------------------------------------------
\section{Handoff to Book IV}
%% ---------------------------------------------------------------

Book~IV begins when the question shifts from \emph{how local certified
structure transfers} to \emph{whether all such transfer data can be
assembled into a single universal object}.  The observational-transfer
arena of Book~III --- objects with stabilized quotients, lawful
transfer maps, and defect ledgers --- is precisely the POT fiber data
that Book~IV will assemble into the universal completion.

The transition from Book~III to Book~IV is:
\begin{quote}
  local admissible observables
  $\;\longrightarrow\;$
  stabilized persistent structure
  $\;\longrightarrow\;$
  lawful transfer across contexts
  $\;\longrightarrow\;$
  universal source completion.
\end{quote}
Book~III supplies the first three terms.  Book~IV supplies the fourth.

%% ---------------------------------------------------------------
\section{Dependency Ledger}
%% ---------------------------------------------------------------

\renewcommand{\arraystretch}{1.4}
\noindent
\begin{tabular}{@{}p{4.0cm}p{0.8cm}p{3.6cm}p{4.6cm}@{}}
\hline
\textbf{Theorem} & \textbf{St.} & \textbf{Depends on} & \textbf{Exports to} \\
\hline
\ref{lem:persistent-quotient-visible}\ Persist.\ obs.\ quotient-vis.
  & [U] & SR~\ref{sr:quotient-visible}, Def.~\ref{def:persistent-obs}
  & \ref{prop:persistent-obs-arena}, \ref{thm:governing} \\

\ref{lem:transfer-preserves-compat}\ Transfer preserves compat.
  & [U] & Def.~\ref{def:transfer-map}
  & \ref{prop:compat-preserved-refinement} \\

\ref{lem:composition-welldef}\ Composition well-def.
  & [U] & Def.~\ref{def:composable-chain}
  & \ref{prop:OT-morphisms}, \ref{prop:OT-arena} \\

\ref{lem:obstructions-quotient-visible}\ Obstructions quot.-vis.
  & [U] & Defs.~\ref{def:persistent-obs}--\ref{def:transport-obstruction}
  & \ref{prop:obstruction-blocks-persistence}, \ref{cor:obstruction-blocks-source} \\

\ref{prop:compatible-comparison-structure}\ Compat.\ comparison
  & [U] & \ref{lem:composition-welldef}, Def.~\ref{def:transfer-compatibility}
  & \ref{cor:compatible-for-source}, \ref{thm:governing} \\

\ref{cor:compatible-for-source}\ Source-ready compat.
  & [U] & \ref{prop:compatible-comparison-structure}
  & Book~IV POT construction \\

\ref{thm:unified-coercive}\ Unified Coercive-Transport Ineq.
  & [U] & Books~I--II, Defs.~\ref{def:stability-functional}--\ref{def:boundary-defect}
  & \ref{thm:balance-composable}, \ref{thm:governing}; all branch books \\

\ref{thm:balance-composable}\ Composable balance
  & [U] & \ref{thm:unified-coercive}, Bk.~I ledger additivity
  & \ref{thm:coercive-inequality}, \ref{thm:governing} \\

\ref{thm:coercive-inequality}\ Finite Transfer Coercive Ineq.
  & [C] & \ref{thm:balance-composable}, coercivity hypotheses
  & Branch books (NS, YM, GR imports) \\

\ref{cor:obstruction-blocks-source}\ Obstruction blocks source
  & [U] & \ref{prop:obstruction-blocks-persistence}, \ref{cor:compatible-for-source}
  & Book~IV; branch books \\

\ref{prop:OT-arena}\ OT comparison arena
  & [U] & \ref{prop:OT-welldef}, \ref{prop:OT-morphisms}
  & Book~IV POT; \ref{thm:governing} \\

\ref{thm:governing}\ Obs.\ Transfer Thm.\ of Spine
  & [U] & All of Bk.~III
  & Books~IV--VII; all branch books \\

\ref{thm:depth-sum}\ Depth-Sum Convergence
  & [U] & Van Hove geometry; Bk.~I finiteness
  & NS Branch Book (weighted boundary bounds) \\

\ref{sch:icdc-abstract}\ ICDC Schema (abstract)
  & [O] & \ref{thm:unified-coercive}; five conditions
  & Branch books Part~III; YM/NS/GR/SCC ICDC \\

\ref{sch:topological-conservation}\ Topological conservation
  & [C] & Gate structure G1--G3
  & GR Branch Book remark \\

\ref{sch:hopf-helicity}\ Hopf helicity
  & [C] & Admissibility of $\Phi$
  & GR Branch Book remark \\
\hline
\end{tabular}
\renewcommand{\arraystretch}{1}

\medskip

\begin{scholium}[\status{R}]\label{sch:book-iii-honest-status}
The Unified Coercive-Transport Inequality
(Theorem~\ref{thm:unified-coercive}) is the load-bearing core of
Book~III and is unconditional: it depends only on the defect ledger of
Book~I and the LS-2 boundary law of Book~II.  The Finite Transfer
Coercive Inequality (Theorem~\ref{thm:coercive-inequality}) is
conditional, requiring extra coercivity hypotheses.  This conditionality
should remain explicit in every downstream citation.  No branch book
may promote the conditional result to unconditional force without
discharging the extra hypotheses within the branch.
\end{scholium}

\end{document}
